\documentclass[../main.tex]{subfiles}

\begin{document}

\chapter{CDCL}

Conflict Driven Clause Learning (CDCL), is \textit{the} modern SAT solving 
algorithm. It builds off of the ideas of DPLL and takes it to a whole new level.
The ``conflict driven'' part of CDCL means that we analyze conflicts (contradictions)
and see what we can learn from them. The ``clause learning'' part is exactly what it 
sounds like, we add new clauses to our database. In other words, CDCL is about 
analyzing why conflicts occur when we solve and using those conflicts to learn 
new clauses that will help enforce new rules as to what values variables can take,
essentially changing the CNF equation into one that's logically equivalent, but 
easier to solve. I won't lie to you, I was not able to implement this algorithm the 
first time I tried. I hope and pray that this time I will be successful, so that 
\textit{we} will be successful together dear reader. 

\textit{Aside: why have we been opting for indices so often? You see, although
we live in the land of pseudocode or machine-agnosticism, implementations don't
(duh). So in our previous chapter in DPLL, there were times where we used clause 
indices instead of pointers to clauses (obviously using copies would be terrible 
for performance and memory), but why did we use indices? Now that we're implementing 
CDCL, we're going to be adding clauses to the database quite often, and that can 
trigger a resize on the database. If we used pointers to clauses, after a resize 
triggers, the memory address of each clause would change, thus making the pointers 
invalid. We would need to add new code after the insertion of each clause to make 
sure a resize didn't trigger, and if one did, we'd need to update each pointer for 
each reason clause and watched literal, which would be very difficult because after 
a resize, we now have no way to tell which clauses mapped to which reasons and watched 
literals, because, well, the old mapping isn't valid and contains no information about 
the new addresses for everything.
}

\footnote{More information about CDCL can be found here: 
    \url{https://www.cis.upenn.edu/~cis1890/files/Lecture4.pdf}}
\section{Clause Learning}

The first step in creating a CDCL solver is to learn and use new clauses. There are 
many versions involving different ways of backtracking, variable assigning, and even 
how we learn new clauses, but for now, we will keep our implementation as close to 
our iterative DPLL version and work our way up to the more advanced and widely used 
versions of the algorithm. 

In order to learn a new clause, we first need to understand the learning process. 
The big idea is that we will encounter conflicts and contradictions while assigning 
and propagating new literals, and we want to learn from those conflicts and contradictions 
in a way that would ideally prevent us from making the same mistake again. Our unit 
propagation method is the perfect tool for this because it already analyzes our clauses 
for us and can detect when a contradiction has occurred. Rather than returning True or 
False to indicate a successful or unsuccessful propagation respectively, when we 
encounter a contradiction, we want to return the clause that is unsatisfied. You may know that
when we propagate, there could be more than one clause that is unsatisfied, but just not 
yet evaluated to be unsatisfied. Even though there could be many conflicting clauses at
once, we just want to return the first one that we find. You may wonder why not learn from 
all of them, and the answer is that by resolving the bad decision/propagation we made, 
we will (hopefully) prevent ourselves from doing more propagations that lead to those 
extra unsatisfied clauses.

After we return the conflicting clause to the main solving method, we then want to 
perform ``clause resolution.'' As stated in the beginning, there are many versions of 
clause resolution, mainly differing in when they/how they stop resolving. For our basic 
CDCL implementation, we will resolve the conflicting clause until it contains no more 
propagated literals, ie, it's a clause of only decisions. The reason why this is a good 
stopping point is because we made all decisions, and because of that, clause learning 
will essentially allow us to learn which decisions we can't make together. 

Lastly, after we resolve on the learned clause, we want to add it to the database,
and we want to initialize its watched literals so that we can use it as if it were 
part of our initial CNF equation. Before we discuss how clause resolution works at 
a theory level, let's first discuss a few implementation changes we'll need to make.

\begin{lstlisting}
    public interface:
        SATSolver(string input)
        bool solveCNF()

    private interface:
        enum var
        {
            False
            True
            Unassigned
        }

        TrailEntry {
            int lit,
            int decisionLevel,
            int clauseIdx
        }

        DecisionLevel {
            int idx,
            bool triedOtherBranch
        }

        void parse(string equation)
        bool propagate(int[] decisionLevel, int decisionLevel)
        bool createWatched()
        bool createWatched(int idx, int decisionLevel)

        int numVars
        int numClauses
        int trailHead

        int[][] clauses
        
        TrailEntry[] trail
        DecisionLevel[] trailStarts
        int[] var_to_trail

        optional var[] vars
        map<int, int[]> var_to_clause
        map<int, (int, int)> clause_to_vars
\end{lstlisting}

The \textit{var\_to\_trail} field is the only change we'll need to make, and it 
is just a mapping from variable index to trail index. If you enforce the invariant 
that each variable maps to exactly one place in the trail and no two variables 
map to the same place in the trail at the same time, you don't need to initialize 
these values to anything. 

The next thing we'll look at is our \textit{propagate()} code since it needs very 
minimal changes. The change we're making now is that instead of returning a boolean
value to indicate whether or not propagation was successful, we're now going to return 
a clause if we encounter a contradiction and nothing otherwise.

\begin{lstlisting}
    bool propagate(int[] assignment, int decisionLevel) {
        while (trailHead < trail.size) {
            int prop = trail[trailHead].lit
            trailHead = trailHead + 1
            ...
            for (i = 0 up to watched_clauses.size) {
                clause_idx = watched_clauses[i]
                int[] clause = clauses[clause_idx]
                (int, int) watches = clause_to_var[clause_idx]
                ...
                if (relocated == False) {
                    ...
                    int other_watch = clause[watches.second];
                    int idx = abs(other_watch)
                    if (assignment[idx] == Unassigned) {
                        ...
                        append(trail, {other_watch, decisionLevel, clause_idx})
                    } else {
                        return clause
                    }
                }
            }
        }
        return null
    }
\end{lstlisting}

And before we look at changes to our solving method, we're going to create a new 
overload for \textit{createWatched()} that takes in the index of a clause. Its 
only purpose is to be called on newly learned clauses, which don't have the same 
guarantees as the other overload. You see, the learned clause should hopefully 
NEVER be unsatisfied when this method operates over it. It should ALWAYS return 
True, so if it ever doesn't, it's a good debugging tool. Here's how it works:
we traverse over the clause of the passed in index, counting the number of 
True, False, and Unassigned literals it has. We will then search for two 
indices that could be watched literals. Since the clause should never be 
unsatisfied, we are guaranteed to always find one (if our implementation is correct).
If we can't find another and our clause has more than one literal, we will initialize 
the second watched literal to any arbitrary place in the clause (with the caveat 
that it can't be the same index as our first watched literal). We will then check if the 
clause is unit, and if it is, we will add it to our trail, assign the Unassigned variable,
etc. Otherwise, we do nothing.

\begin{lstlisting}
    bool createWatched(int index, int decisionLevel) {
        clause = clauses[index]
        if (clause.size == 0) {
            return False
        }

        bool foundFirst = False 
        bool foundSecond = False 
        int firstIdx = -1
        int secondIdx = -1

        int numUnassigned = 0
        int numFalse = 0
        int numTrue = 0

        for (i = 1 up to clause.size) {
            int idx = abs(clause[i])
            if (assignment[idx] == Unassigned) {
                numUnassigned = numUnassigned + 1
                if (foundFirst == False) {
                    foundFirst = True
                    firstIdx = i
                } else if (foundSecond == False) {
                    foundSecond = True 
                    secondIdx = i
                }
            } else if (evals_false(clause[i])) {
                numFalse = numFalse + 1
            } else {
                numTrue = numTrue + 1

                if (foundFirst == False) {
                    foundFirst = True
                    firstIdx = i
                } else if (foundSecond == False) {
                    foundSecond = True 
                    secondIdx = i
                }
            }
        }

        if (clause.size >= 2) {
            if (foundSecond == False) {
                secondIdx = (firstIdx + 1) % clause.size
            }

            append(var_to_clause[clause[firstIdx]], index)
            append(var_to_clause[clause[secondIdx]], index)
            
            clause_to_var[index] = {firstIdx, secondIdx}
        } else {
            append(var_to_clause[clause[firstIdx]])
            clause_to_var[index] = {firstIdx, firstIdx}
        }

        if (numTrue == 0 && numUnassigned == 1) {
            append(trail, {clause[firstIdx], decisionLevel, index})
            var_to_trail[abs(clause[firstIdx])] = trail.size() - 1

            int var = clause[firstIdx]
            int idx = abs(var)

            if (var > 0) {
                assignment[idx] = True
            } else {
                assignment[idx] = False
            }
        }

        return True
    }
\end{lstlisting}

Because we don't know anything about this clause, we need to gather as much 
information as we can right away before we decide how watched literals are 
spread out on the clause. When the clause has more than one literal, we 
use ``secondIdx = (firstIdx + 1) \% clause.size" because that prevents us 
from adding a watched literal past the end of the clause, if \textit{firstIdx}
is already on the edge.

The last method we need to discuss is our \textit{solve} loop. Clause 
resolution happens after backtracking because it relies on the current 
trail to understand
\footnote{It follows a mathematical process to transform the conflict clause 
into one that prevents that same conflict from ever occurring again. The 
algorithm does not ``understand'' anything. Do not anthropomorphize computer 
algorithms, especially artificial intelligence.}
why a conflict occurred. If we were to undo propagations, then the resolution 
operator would have no reference as to what ``most recently'' means, and 
therefore, clause resolution wouldn't work.

\begin{lstlisting}
    bool solve() {
        var[1..numVars] assignment

        for each (var in assignment) {
            var = Unassigned
        }

        if (createWatched(assignment) == False) {
            return False
        }

        int decisionLevel = 0

        if (trail.size > 0) {
            append(trailStarts, {0, true})

            if (propagate(assignment, decisionLevel) == False) {
                return False
            }
        }

        while (True) {
            if (trailHead == trail.size) {
                //Assignment code
                int firstUnassigned = -1
                for (i = 1 up to numVars) {
                    if (assignment[i] == Unassigned) {
                        firstUnassigned = i
                        break
                    }
                }

                if (firstUnassigned == -1) {
                    return True
                }

                assignment[firstUnassigned] = True
                append(trail, {firstUnassigned, -1})
                append(trailStarts, {trail.size - 1, False})
                trailHead = trailStarts.last.idx
            }

            optional int[] contradiction = propagate(assignment)

            if (contradiction has a value) {
                int[] conflict = contradiction.value

                bool propagated_found = True
                while (True) {
                    propagated_found = False
                    int idx = 0

                    for (i = 1 up to conflict.size) {
                        int var = conflict[i]
                        int v_idx = abs(var)

                        int t_idx = var_to_trail[v_idx]
                        if (trail[t_idx].reasonIdx has a value) {
                            propagated_found = True
                            idx = max(idx, t_idx)
                        }
                    }

                    if (propagated_found == False) {
                        break
                    }

                    int lit = trail[idx].lit
                    int[] reason = clauses[trail[idx].reasonIdx]

                    append(conflict, reason)

                    for (i = 0 up to conflict.size) {
                        if (abs(conflict[i]) == abs(lit)) {
                            swap(conflict[i], conflict.last)
                            erase(conflict, conflict.last)
                        } else {
                            i = i + 1
                        }
                    }

                    set<int> duplicates
                    for (i = 0 up to conflict.size) {
                        if (contains(duplicates, conflict[i])) {
                            swap(conflict[i], conflict.last)
                            erase(conflict, conflict.last)
                        } else {
                            insert(duplicates, conflict[i])
                            i = i + 1
                        }
                    }
                }

                while (trailStarts.size > 0) {
                    DecisionLevel decision = trailStarts.last

                    while (trail.size > decision.idx + 1) {
                        int lit = trail.last.lit
                        assignment[abs(lit)] = Unassigned
                        var_to_trail[abs(lit)] = -1
                        erase(trail, trail.last)
                    }

                    if (decision.triedFalse == False) {
                        TrailEntry entry = trail.last

                        entry.lit = -entry.lit
                        assignment[abs(entry.lit)] = False
                        var_to_trail[abs(entry.lit)] = -1
                        decision.triedFalse = True

                        trailHead = decision.idx
                        break
                    } else {
                        int lit = trail.last.lit
                        assignment[abs(lit)] = Unassigned
                        var_to_trail[abs(lit)] = -1
                        erase(trail.last)
                        erase(trailStarts.last)
                    }
                }

                if (trailStarts.size == 0) {
                    return False
                }

                append(clauses, conflict)
                createWatched(clauses.size)

                continue
            }
        }
    }
\end{lstlisting}

I want to point out that this isn't nevessarily CDCL yet. This is more like DPLL, but with 
adding clauses to our database. The reason as to why this is, is because whenever we recover 
from a contradiction, we only go back one level. You see, CDCL does away with the idea of a
``tried other branch" flag and relies on the learned clauses to guide what variable assignments
are possible and which ones are contradictory. However, because we only backtrack one level, there 
is nothing to actually force new variable assignments in cases where the learned clause isn't a unit 
clause. It can happen, but it doesn't always happen, so we still need to try the other assignment 
since for the most part, the learned clause isn't used.

\subsection{Clause Resolution}

Although we've already seen the pseudocode for clause resolution, it's not discussed very 
much, and I want us to understand how it works better. As we make decisions and propagations, 
we form this thing called an \textbf{implication graph}. In a DPLL style solver, our assignments 
and propagations form a directed acyclic graph (I will not be proving this), and by traversing the
implication graph, we can find out what decisions led to the contradiction.

To learn more about implication graphs, let's look at one and discuss what it means.
We will use the implication graphs from here (Figures 4.1 and 4.2): 
\begin{center}
    \small{\url{https://www3.cs.stonybrook.edu/~cram/cse505/Fall20/Resources/cdcl.pdf}}
\end{center}

We will start with Figure 4.1:

\fbox{
    \begin{tikzpicture}[
        >=Stealth,
        every node/.style={font=\large},
        implication/.style={->, thick}
    ]

    % Nodes
    \node (x31) {$x_{31}=0\ @\ 3$};
    \node (x2)  [below right=0.8cm and 1.8cm of x31] {$x_2=0\ @\ 5$};
    \node (x1)  [below left=0.8cm and 1.8cm of x2] {$x_1=0\ @\ 5$};
    \node (x3)  [below right=0.8cm and 1.8cm of x1] {$x_3=0\ @\ 5$};
    \node (x4)  [above right=0.8cm and 1.8cm of x3] {$x_4=1\ @\ 5$};

    % Edges
    \draw[implication] (x31) -- node[above left] {$\omega_1$} (x2);
    \draw[implication] (x1)  -- node[above left] {$\omega_1$} (x2);
    \draw[implication] (x1)  -- node[below left] {$\omega_2$} (x3);
    \draw[implication] (x2)  -- node[above right] {$\omega_3$} (x4);
    \draw[implication] (x3)  -- node[below right] {$\omega_3$} (x4);

    \end{tikzpicture}
}

It corresponds to the following CNF equation: 
$$ (\omega_1) \wedge (\omega_2) \wedge (\omega_3) =
(x_1 \vee x_{31} \vee \neg x_2) \wedge (x_1 \vee \neg x_3) \wedge (x_2 \vee x_3 \vee x_4)$$

Each node of the graph follows this form: $x_i = v @ d$, but what does it mean?

\begin{center}
    This is the notation for assignment of a variable at a certain decision level:
    $$x_i = v@d$$
    where $i$ denotes the index of the variable being assigned, $v \in \{0, 1\}$
    (alternatively, $v \in \{F, T\}$), and $d$ is the decision level for which that 
    assignment was executed.
\end{center}

So in our graph, $x_{31} = 0 @ 3$ means variable 31 was assigned 0 (or False) at
decision level 3. What about the edges? The edges are clauses, and they connect
variables to each other, but what's the relationship? Let's go left to right in our 
implication graph. Variables 31 and 1 are assigned False at decision levels 3 and 5
respectively, and through clause 1, they both have an edge to variable 2, which is also 
assigned False at decision level 5. If we look at the original CNF equation, we see that 
in clause 1, $x_1$ and $x_{31}$ are not negated, so they are False literals. That leaves 
$x_2$ Unassigned, so we propagate in order to satisfy the clause. We assign $x_2 = F$ 
because it's negated, so its literal will evaluate to True. Likewise, $x_1$ has another 
edge of clause 2 ($\omega_2$) to $x_3$. If we look at clause 2, we see that $x_1 = F$ 
forces $\neg x_3 = T \equiv x_3 = F$. Lastly, we see that $x_2$ and $x_3$ have an edge 
to $x_4$ through $\omega_3$, forcing its assignment to be $x_4 = T$.

That was pretty easy to understand, so now let's look at an implication graph with 
contradictions. From Figure 4.2:

\noindent
\fbox{
\resizebox{\textwidth}{!}{%
\begin{tikzpicture}[
    >=Stealth,
    every node/.style={font=\large},
    implication/.style={->, thick}
]

% Left side
\node (x31) {$x_{31}=0\ @\ 3$};

\node (x2)
    [below right=0.8cm and 1.8cm of x31]
    {$x_2=0\ @\ 5$};

\node (x1)
    [below left=0.8cm and 1.8cm of x2]
    {$x_1=0\ @\ 5$};

\node (x3)
    [below right=0.8cm and 1.8cm of x1]
    {$x_3=0\ @\ 5$};

\node (x4)
    [above right=0.8cm and 1.8cm of x3]
    {$x_4=1\ @\ 5$};

% Right side
\node (x5)
    [above right=0.8cm and 2.5cm of x4]
    {$x_5=0\ @\ 5$};

\node (x6)
    [below right=0.8cm and 2.5cm of x4]
    {$x_6=0\ @\ 5$};

\node (x21)
    [below left=1.0cm and 1.8cm of x6]
    {$x_{21}=0\ @\ 2$};

\node (kappa)
    [right=2.5cm of x6]
    {$\kappa$};

% Left implication graph
\draw[implication]
    (x31) -- node[above left] {$\omega_1$} (x2);

\draw[implication]
    (x1) -- node[above left] {$\omega_1$} (x2);

\draw[implication]
    (x1) -- node[below left] {$\omega_2$} (x3);

\draw[implication]
    (x2) -- node[above right] {$\omega_3$} (x4);

\draw[implication]
    (x3) -- node[below right] {$\omega_3$} (x4);

% Right implication graph
\draw[implication]
    (x4) -- node[above left] {$\omega_4$} (x5);

\draw[implication]
    (x4) -- node[below left] {$\omega_5$} (x6);

\draw[implication]
    (x21) -- node[below right] {$\omega_5$} (x6);

\draw[implication]
    (x5) -- node[above right] {$\omega_6$} (kappa);

\draw[implication]
    (x6) -- node[below right] {$\omega_6$} (kappa);

\end{tikzpicture}
}
}

\noindent
This graph corresponds to the CNF equation: 
$$ (\omega_1) \wedge (\omega_2) \wedge (\omega_3) \wedge (\omega_4) \wedge (\omega_5) \wedge (\omega_6) 
=
(x_1 \vee x_{31} \vee \neg x_2)_1 \wedge (x_1 \vee \neg x_3)_2 \wedge (x_2 \vee x_3 \vee x_4)_3
$$
$$
\wedge (\neg x_4 \vee \neg x_5)_4 \wedge (x_{21} \vee \neg x_4 \vee \neg x_6)_5 \wedge (x_5 \vee x_6)_6
$$

Up until $x_4 = 1 @ 5$, this graph is the same as the one from above. After that,
we have some more decisions and propagations, that I bet you can read by now, up until 
we reach this symbol: $\kappa$, which has two edges via $\omega_6$. This node, $\kappa$,
corresponds to a contradiction within clause $\omega_6$ because its literals are all False.

\noindent
When we resolve this conflict, we first start with the clause 
$$(x_5 \vee x_6)$$
as that is the clause in which our conflict occurred. Next, we find the most recently propagated literal, 
which we will assume to be $x_6$ (the most recently propagated literal depends on our implementation 
of propagation, so for simplicity's sake, we will assume numerical order to settle ties in 
implication orderings on the same level). $x_6$ was propagated because of the clause
$$(x_{21} \vee \neg x_4 \vee \neg x_6)$$
Combining this with the current conflict clause yields:
$$(x_5 \vee x_6 \vee x_{21} \vee \neg x_4 \vee \neg x_6)$$
which we process to remove all instances of the conflicting variable, leading to a resolution 
clause of: 
$$(x_5 \vee x_{21} \vee \neg x_4)$$
In this clause, the most recently propagated literal is $x_5$, so we grab its reason clause ($\omega_4$)
and combine once again to get: 
$$(\neg x_4 \vee \neg x_5 \vee x_5 \vee x_{21} \vee \neg x_4)$$

\noindent
Simplifying that clause gives us: $$(\neg x_4 \vee x_{21})$$
The only propagated literal in the 
clause is $x_4$, so we grab its reason clause ($\omega_3$) and combine again into:
$$(x_2 \vee x_3 \vee x_4 \vee \neg x_4 \vee x_{21})$$
This simplifies down into:
$$(x_2 \vee x_3 \vee x_{21})$$
Our most recently propagated literal is now $x_3$. Repeating 
the process again, we get an unsimplified resolution clause of: 
$$(x_2 \vee x_3 \vee x_{21} \vee x_1 \vee \neg x_3)$$ 
Simplifying it down, we now get:
$$(x_2 \vee x_{21} \vee x_1)$$
Our last propagated literal in the clause is $x_2$. We 
combine our incomplete resolution clause with $\omega_1$ to get:
$$(x_2 \vee x_{21} \vee x_1 \vee x_1 \vee x_{31} \vee \neg x_2)$$
This transforms into 
$$(x_1 \vee x_{21} \vee x_{31})$$
which contains only decision literals.

What does this new clause, which we'll call $\omega_7$, tell us? It tells us that 
all three decision literals \textbf{cannot} all be False (ie, at least one of them 
\textbf{must} be True).

\subsection{Clause Learning - Results}

Here are the results for Clause Learning:

\noindent
uf20.91: \textbf{425ms}, up 98ms from Iterative DPLL 

\noindent
UUF50.218.1000: \textbf{9652ms}, up 3287ms from Iterative DPLL

Why did our runtime go up? Well, clause learning isn't easy, and the more 
contradictions we encounter, the more expensive clause learning we're performing. 
Another reason is that we're increasing the memory load every time we add in a new
clause. Luckily, our testing sets are small, but on large problems, this would be a 
much larger issue. So then why do modern solvers use clause learning? Because there 
are ways to make it much faster. The first step would be to implement non-chronological 
backjumping, where instead of moving up just one decision level in our trail, we move up 
multiple at a time. The second step would be to implement First UIPs, which makes clause 
learning a less expensive operation (doesn't fix the memory load though), and it
combined with non chronological backjumping allows us to instantly and fully use our 
learned clause to its potential. Although we have taken a small step back, we will 
see what leaps forward we will make in the next sections.

\textit{Aside: Why did we get rid of pure literal elimination? The problem with pure literal 
elimination is that once we start learning clauses, previously pure literals may no longer 
be pure anymore once we add in a new clause. Because pure literals aren't born from decisions/
assignments, they aren't as ``concrete'' in their status. (Ie, propagation \textbf{forces} new 
literal assignments, pure literal elimination is more of a ``rule of thumb''). If we had a pure 
literal, say $5$, and we added in a clause containing $-5$, unfortunately, $5$ is no longer pure, 
and anything propagated because of $5$ also would need to be undone because the reason why they 
were propagated/assigned those values in the first place doesn't hold anymore. Because we saw that 
pure literal elimination wasn't that helpful to begin with and with all these reasons as to why it 
wouldn't easily work under CDCL, it would just be much better to not have it at all.}

\begin{FVerbatim}
On all 1000 equations in uf20-91:

Clause Learning:           425ms (~0.43 seconds).
Iterative DPLL (AbP):      328ms (~0.33 seconds).
Iterative DPLL (PbA):      324ms (~0.32 seconds).
Watched Literals:          349ms (~0.35 seconds).
Pure Literal Elimination:  1548ms (~1.55 seconds).
Unit Propagation:          868ms (~0.87 seconds).
Recursive Backtracking:    1,540,613ms (~25.68 minutes).
Brute Force:               1,770,851ms (~29.51 minutes).
Short-Circuited (Basic):   49,874ms (~50 seconds).
Parallel:                  10,589ms (~10.5 seconds).

On all 1000 equations in UUF50.218.1000:

Clause Learning:          9,652ms (~9.65 seconds).
Iterative DPLL (AbP):     6,369ms (~6.37 seconds).
Iterative DPLL (PbA):     6,365ms (~6.37 seconds).
Watched Literals:         7,130ms (~7.13 seconds).
Pure Literal Elimination: 231,107ms (~3.85 minutes).
Unit Propagation:         140,041ms (~2.34 minutes).
\end{FVerbatim}

\section{Non-Chronological Backjumping}

An issue I mentioned earlier with our initial clause learning algorithm is that we backtrack 
only one level per contradiction. Not one level at a time, just one level. The reason why this 
is an issue is that we generally never get to use what we learned from conflict resolution to 
flip certain assignments, which is why we still relied on the DPLL notion of ``trying the other 
branch.'' Non-chronological backtracking, or ``backjumping'' is the new technique we'll use 
to fully exploit the learned clause. 

We only need to make a very minimal interface change. Now that we won't use a ``triedFalse''
variable for our trail starts, our trail starts array can now just be an integer array 
where each element is just the index in the trail where a new decision level starts.
\begin{lstlisting}
    public interface:
        SATSolver(string input)
        bool solveCNF()

    private interface:
        ...
        int[] trailStarts
        ...
\end{lstlisting}

Next, since we're only changing the backtracking logic, we can keep 
everything else the same in our algorithm. Before we move onto code, answer this:
what decision level should be backjump to after we learn a clause? Suppose we learn 
this clause: 

$$ (1 \vee 2 \vee \neg 4) $$

And each literal's decision level is its index (1 was assigned on decision level 1, 
2 was assigned on decision level 2, 4 was assigned on level 4). What decision level 
should be backjump to? Should we backjump to 1? 2? Stay on 4? We want to backjump to 
decision level 2 because that way, every literal in the clause other than $4$ will 
be False, and $4$ will be Unassigned. That way, we can propagate immediately on the 
learned clause! So the rule for backjumping after learning a new clause is to backjump 
to the second largest decision level present in the clause. Here's how we'll do it:

\begin{lstlisting}
bool solve() {
    var[1..numVars] assignment

    for each (var in assignment) {
        var = Unassigned
    }

    if (createWatched(assignment) == False) {
        return False
    }

    int decisionLevel = 0

    if (trail.size > 0) {
        append(trailStarts, {0, true})

        if (propagate(assignment, decisionLevel) == False) {
            return False
        }
    }

    while (True) {
        if (trailHead == trail.size) {
            //Perform an assignment
            ...
        }

        optional int[] contradiction = propagate(assignment)

        if (contradiction has a value) {
            //resolve conflict
            ...

            if (conflict.size == 0) {
                return False
            }

            int maxLevel = decisionLevel
            int secondLargest = -1
            for (i = 0 up to conflict.size) {
                int idx = abs(conflict[i])
                int decisionLvl = trail[var_to_trail[idx]].decisionLevel;
                if (decisionLvl != maxLevel) {
                    secondLargest = max(secondLargest, decisionLvl)
                }
            }

            int backjumpTo = secondLargest

            while (trail.size > 0 && trail.last.decisionLevel > backjumpTo) {
                int decision_start = trailStarts.last
                while (trail.size > decision_start) {
                    int lit = trail.last.lit
                    int idx = abs(lit)
                    assignment[idx] = Unassigned
                    var_to_trail[idx] = -1
                    erase(trail, trail.last)
                }
                erase(trailStarts, trailStarts.last)
            }

            if (trailStarts.size == 0) {
                return False
            }

            decisionLevel = backjumpTo

            trailHead = trailStarts.last

            append(clauses, conflict)
            createWatched(clauses.size)

            continue
        }
    }
}
\end{lstlisting}

The code remains largely the same. However, we now want to check that the learned clause isn't empty
so we can avoid doing large amounts of work. When can the learned clause become empty? To answer 
that question, let's answer this other question: ``what does an empty clause mean?'' From our 
section on Watched Literals, an empty clause means that the equation is unsatisfiable because 
no variable assignment can satisfy that clause. So if we learn an empty clause, that means that 
we're now learning that there is no satisfactory assignment of variables, thus our equation is 
unsatisfiable.

\subsection{Non-Chronological Backjumping - Results}
Okay so I know I said that we would make leaps forward in the next sections.
I may have lied. Here are the results for Non-Chronological Backjumping:

\noindent
uf20.91: \textbf{380ms}, down 45ms from Clause Learning.

\noindent
UUF50.218.1000: \textbf{11291ms}, up 1639ms from Clause Learning.

Why did our time to solve unsatisfiable equations go up so significantly? It has to do 
with the fact of how we're learning. We're not trying every assignment like brute force, but 
we are slowly learning that no matter what the assignment is, our equation is unsatisfiable. 
This takes a lot of backtracking and memory to learn, possibly attributing to our time increase.
It's hard to say for certain.

\begin{FVerbatim}
On all 1000 equations in uf20-91:

Backjumping:               380ms (0.38 seconds).
Clause Learning:           425ms (~0.43 seconds).
Iterative DPLL (AbP):      328ms (~0.33 seconds).
Iterative DPLL (PbA):      324ms (~0.32 seconds).
Watched Literals:          349ms (~0.35 seconds).
Pure Literal Elimination:  1548ms (~1.55 seconds).
Unit Propagation:          868ms (~0.87 seconds).
Recursive Backtracking:    1,540,613ms (~25.68 minutes).
Brute Force:               1,770,851ms (~29.51 minutes).
Short-Circuited (Basic):   49,874ms (~50 seconds).
Parallel:                  10,589ms (~10.5 seconds).

On all 1000 equations in UUF50.218.1000:

Backjumping:              11,291ms (~11.29 seconds).
Clause Learning:          9,652ms (~9.65 seconds).
Iterative DPLL (AbP):     6,369ms (~6.37 seconds).
Iterative DPLL (PbA):     6,365ms (~6.37 seconds).
Watched Literals:         7,130ms (~7.13 seconds).
Pure Literal Elimination: 231,107ms (~3.85 minutes).
Unit Propagation:         140,041ms (~2.34 minutes).
\end{FVerbatim}

\section{First UIPs}

Our next addition to CDCL is First UIPs. UIP stands for "unit implication point," which 
is called a ``dominator'' in our implication graph. What is a dominator? A node $d$ is said 
to \textit{dominate} another node $n$ if every path in the graph from the entry node to $n$ 
passes through $d$. Think of a binary tree. Our entry node is the the head of the tree, and 
it dominates its two children. Likewise, each of those children dominates their own children 
nodes. So in other words, the UIP in our implication graph is kind of like the reason why 
we had a contradiction in the first place!

Given that summary of what a UIP is, wouldn't the first UIP be the decision we made on 
the current decision level? After all, we wouldn't have had a contradiction if we didn't 
make that decision. Well not necessarily, because in our implication graph, every decision 
we make only has outward edges, whereas all propagated nodes have inward edges as well. 
So while it is true that a first UIP could be a decision variable, it's not true that it's 
\textit{always} a decision variable. Take this previous implication graph as an example:

\noindent
\fbox{
\resizebox{\textwidth}{!}{%
\begin{tikzpicture}[
    >=Stealth,
    every node/.style={font=\large},
    implication/.style={->, thick}
]

% Left side
\node (x31) {$x_{31}=0\ @\ 3$};

\node (x2)
    [below right=0.8cm and 1.8cm of x31]
    {$x_2=0\ @\ 5$};

\node (x1)
    [below left=0.8cm and 1.8cm of x2]
    {$x_1=0\ @\ 5$};

\node (x3)
    [below right=0.8cm and 1.8cm of x1]
    {$x_3=0\ @\ 5$};

\node (x4)
    [above right=0.8cm and 1.8cm of x3]
    {$x_4=1\ @\ 5$};

% Right side
\node (x5)
    [above right=0.8cm and 2.5cm of x4]
    {$x_5=0\ @\ 5$};

\node (x6)
    [below right=0.8cm and 2.5cm of x4]
    {$x_6=0\ @\ 5$};

\node (x21)
    [below left=1.0cm and 1.8cm of x6]
    {$x_{21}=0\ @\ 2$};

\node (kappa)
    [right=2.5cm of x6]
    {$\kappa$};

% Left implication graph
\draw[implication]
    (x31) -- node[above left] {$\omega_1$} (x2);

\draw[implication]
    (x1) -- node[above left] {$\omega_1$} (x2);

\draw[implication]
    (x1) -- node[below left] {$\omega_2$} (x3);

\draw[implication]
    (x2) -- node[above right] {$\omega_3$} (x4);

\draw[implication]
    (x3) -- node[below right] {$\omega_3$} (x4);

% Right implication graph
\draw[implication]
    (x4) -- node[above left] {$\omega_4$} (x5);

\draw[implication]
    (x4) -- node[below left] {$\omega_5$} (x6);

\draw[implication]
    (x21) -- node[below right] {$\omega_5$} (x6);

\draw[implication]
    (x5) -- node[above right] {$\omega_6$} (kappa);

\draw[implication]
    (x6) -- node[below right] {$\omega_6$} (kappa);

\end{tikzpicture}
}
}
$x_4$ is a UIP in our implication graph because every path to $x_5$ and $x_6$ all 
pass through $x_4$'s node.

So what does the first UIP offer us? Resolving clauses up until we encounter the first UIP 
maximizes the level we backjump to and it makes the clauses we learn shorter, not just in length,
but also in how long it takes to resolve a clause. More details can be found here:
\begin{center}
    \small{\url{https://www3.cs.stonybrook.edu/~cram/cse505/Fall20/Resources/cdcl.pdf}}
\end{center}

How do we find the first UIP? The first UIP in a conflict is found when a the number 
of literals in the clause for the \textit{current} decision level is 1. So if our current 
decision level is 6, we resolve only on literals from decision level 6 and stop once the 
clause contains only 1 literal from decision level 6.

\subsection{Implementing First UIPs}

Once again, we only need to touch the conflict resolution code:

\begin{lstlisting}
bool solve() {
    //setup
    ...

    while (True) {
        if (trailHead == trail.size) {
            //Perform an assignment
            ...
        }

        optional int[] contradiction = propagate(assignment)

        if (contradiction has a value) {
            int curLevel = decisionLevel
            int[] conflict = contradiction.value
            int end = trail.size
            int begin = 0
            int resolveOn = 0

            for (i = 0 up to conflict.size) {
                int idx = abs(conflict[i])
                int trail_idx = var_to_trail[idx]

                if (trail[trail_idx].decisionLevel == curLevel) {
                    end = min(end, trail_idx)
                    begin = max(begin, trail_idx)
                    if (trail[trail_idx].reasonIdx is not optional) {
                        resolveOn = trail[trail_idx].lit
                    }
                }
            }

            while (begin != end) {
                int trail_idx = var_to_trail[abs(resolveOn)]
                int[] reason = clauses[trail[trail_idx].reasonIdx]
                append(conflict, reason)
                for (i = 0 up to conflict.size) {
                    if (abs(conflict[i]) == abs(resolveOn)) {
                        erase(conflict, conflict[i])
                    } else {
                        i = i + 1
                    }
                }

                set<int> duplicates;
                for (i = 0 up to conflict.size) {
                    if (contains(duplicates, conflict[i])) {
                        erase(conflict, conflict[i])
                    } else {
                        duplicates.insert(conflict[i]);
                        i = i + 1
                    }
                }

                end = INTEGER_MAX;
                begin = 0;

                for (i = 0 up to conflict.size) {
                    int var = conflict[i];
                    int idx = abs(var)
                    int trail_idx = var_to_trail[idx];

                    if (trail[trail_idx].decisionLevel 
                        == decisionLevel)
                    {
                        end = min(end, trail_idx);
                        begin = max(begin, trail_idx);

                        if (trail[trail_idx].reasonIdx 
                            is not optional)
                        {
                            resolve_var = trail[trail_idx].lit;
                        }
                    }
                }
            } 

            //backjump
            ...

            if (trailStarts.size == 0) {
                return False
            }

            decisionLevel = backjumpTo

            trailHead = trailStarts.last

            append(clauses, conflict)
            createWatched(clauses.size)

            continue
        }
    }
}
\end{lstlisting}

This is not a very efficient way of computing the first UIP (which you'll see in the 
results). However, let's analyze it anyways. It's the normal resolution code, but we 
do something different, which is stop when there's only one literal from the current 
decision level left in the clause. In our pseudocode, we don't compute the counts of 
literals from the current decision level, we have the condition $begin != end$. Why is 
that equivalent? Well \textit{end} is the index of the least recently propagated literal 
in the resolution clause, and \textit{begin} is the index of the most recently propagated 
literal in the resolution clause. So when the most recently and least recently propagated 
literals are the same, there must be only one left, therefore we can stop. Although it's 
clever, it's not efficient, and therefore, not very clever.

\subsection{First UIPs - Results}

Unfortunately our runtime increases again. At this point, I'm willing to accept
that perhaps I was just really really good at implementing DPLL, but since I'm not 
good at implementing CDCL (I literally failed the first time), my implementation 
is not up to par. In the next section, I will identify possible improvements that 
can be made, and hopefully they'll make a \textit{positive} difference.

\noindent
uf20.91: \textbf{406ms}

\noindent
UUF50.218.1000: \textbf{14012ms}

\begin{FVerbatim}
On all 1000 equations in uf20-91:

First UIPs:                406ms (~0.41 seconds)
Backjumping:               380ms (~0.38 seconds).
Clause Learning:           425ms (~0.43 seconds).
Iterative DPLL (AbP):      328ms (~0.33 seconds).
Iterative DPLL (PbA):      324ms (~0.32 seconds).
Watched Literals:          349ms (~0.35 seconds).
Pure Literal Elimination:  1548ms (~1.55 seconds).
Unit Propagation:          868ms (~0.87 seconds).
Recursive Backtracking:    1,540,613ms (~25.68 minutes).
Brute Force:               1,770,851ms (~29.51 minutes).
Short-Circuited (Basic):   49,874ms (~50 seconds).
Parallel:                  10,589ms (~10.5 seconds).

On all 1000 equations in UUF50.218.1000:

First UIPs:               14,012ms (~14.01 seconds).
Backjumping:              11,291ms (~11.29 seconds).
Clause Learning:          9,652ms (~9.65 seconds).
Iterative DPLL (AbP):     6,369ms (~6.37 seconds).
Iterative DPLL (PbA):     6,365ms (~6.37 seconds).
Watched Literals:         7,130ms (~7.13 seconds).
Pure Literal Elimination: 231,107ms (~3.85 minutes).
Unit Propagation:         140,041ms (~2.34 minutes).
\end{FVerbatim}

\section{A Second Round of Improvements}

Ever since we started CDCL, I have been promising massive speedups over DPLL 
and yet our SAT solver has gotten slower each iteration. Before you treat me 
like Jesus and crucify me,
\footnote{My friends told me that it was too funny not to include so it stays.}
I knew and deep down you knew that the way we were implementing was not optimal
to begin with. So in a small way, this is your fault too. 

\subsection{Removing Hashing}

The first thing we're going to do is remove the maps and instead use arrays 
instead so we don't have to deal with hashing. Our interface becomes like so:

\begin{lstlisting}
    public interface:
        SATSolver(string input)
        bool solveCNF()

    private interface:
        ...
        int[][] var_to_clause
        (int, int)[] clause_to_var
\end{lstlisting}

In case you need a reminder as it has been a while, these are the data structures 
used to make watched literals happen. Before they were maps where the keys were variables 
or clause indices, but now, each index in the arrays are the keys and the values at those 
indices are, well, the values. There is a small implementation detail concerning 
\textit{clause\_to\_var}, which is that now that the number of clauses are not fixed, 
everytime we add in a new clause, we need to add in a new entry to \textit{clause\_to\_var}
before we let our \textit{createWatched(int, int)} overload handle assigning the values 
to that entry. Basically, to avoid going out of bounds for the data structure, we need to 
change the bounds before we access the new boundary. While we could let 
\textit{createWatched(int, int)} handle the allocation as well, it would make it less 
versatile as it would have to only operate on the last clause in the database instead of 
any clause in the database.

The second point to mention is that now our propagation logic doesn't work because we 
need to create a mapping from literals (positive and negative) to indices (nonnegative).
We use the following mapping:
$$ 2 * abs(lit) + (lit < 0) $$
This is how the mapping works out (left column is literal value, right column is index):
\begin{FVerbatim}
 1 → 0
-1 → 1
 2 → 2
-2 → 3
 3 → 4
-3 → 5
  ...
\end{FVerbatim}

The last step is to remove the duplicates set we use in clause resolution.
We're going to add in an array of booleans to our SAT solver where index $i$ 
is whether or not we've seen that variable during our clause resolution.
However, since our (my) implementation of first UIP clause resolution is 
such a big fixer-upper, it deserves its own subsection.

\subsection{Fixing Clause Resolution}

Here is our current clause resolution code:
\begin{lstlisting}
if (contradiction has a value) {
    int curLevel = decisionLevel
    int[] conflict = contradiction.value
    int end = trail.size
    int begin = 0
    int resolveOn = 0

    for (i = 0 up to conflict.size) {
        int idx = abs(conflict[i])
        int trail_idx = var_to_trail[idx]

        if (trail[trail_idx].decisionLevel == curLevel) {
            end = min(end, trail_idx)
            begin = max(begin, trail_idx)
            if (trail[trail_idx].reasonIdx is not optional) {
                resolveOn = trail[trail_idx].lit
            }
        }
    }

    while (begin != end) {
        int trail_idx = var_to_trail[abs(resolveOn)]
        int[] reason = clauses[trail[trail_idx].reasonIdx]
        append(conflict, reason)
        for (i = 0 up to conflict.size) {
            if (abs(conflict[i]) == abs(resolveOn)) {
                erase(conflict, conflict[i])
            } else {
                i = i + 1
            }
        }

        set<int> duplicates;
        for (i = 0 up to conflict.size) {
            if (contains(duplicates, conflict[i])) {
                erase(conflict, conflict[i])
            } else {
                duplicates.insert(conflict[i]);
                i = i + 1
            }
        }

        end = INTEGER_MAX;
        begin = 0;

        for (i = 0 up to conflict.size) {
            int var = conflict[i];
            int idx = abs(var)
            int trail_idx = var_to_trail[idx];

            if (trail[trail_idx].decisionLevel 
                == decisionLevel)
            {
                end = min(end, trail_idx);
                begin = max(begin, trail_idx);

                if (trail[trail_idx].reasonIdx 
                    is not optional)
                {
                    resolve_var = trail[trail_idx].lit;
                }
            }
        }
    } 
}
\end{lstlisting}
\noindent
It's pretty long and inefficient. This was written before we understood implication
graphs fully. Now we do understand implication graphs, so let's exploit that structure.

The idea is this: instead of rescanning over and over to make sure our resolution clause
has just one current level variable, we'll keep track of which variables we've seen as 
well as which ones are from this decision level. Then, as we perform clause resolution, 
we'll update the seen and count variables to reflect the current state of the resolution 
clause, then once we've encountered our stopping condition, we'll use which variables 
we've kept marked as seen to construct our learned clause directly. This is very different 
from the way we've been performing clause resolution, which involved repeatedly scanning the 
resolution clause and recomputing how far along we are in the UIP process. 

How are we using the implication graph in our new method? Well since the implication 
graph is an acyclic graph, that means we can use its structure for our UIP traversal. 
There is a relationship between the implication graph and our trail, in that our trail 
is just a compressed version of the implication graph, meaning it doesn't keep track of 
which variable assignments led to which propagations. However, they both follow the same 
structure in that the edge of our trail are the edges of our implication graphs and the 
order of variables in the trail and graphs are the same. Because the implication graph 
is acyclic, once we encounter a variable during clause resolution, we will never encounter 
it again while we're resolving. Thanks to that property, we can traverse the trail backwards 
without having to worry about going over the same parts of the trail/graph, which was 
something that we were sort of doing in the old method when we rescanned the whole 
conflict clause, whereas now, we are only concerned with the differences to the resolution
clause made across each iteration. 

Our resolution algorithm becomes as follows:
\begin{lstlisting}
if (contradiction has a value) {
    int[] conflict = contradiction.value
    int levelCount = 0
    int curLevel = decisionLevel

    for each (int var in conflict) {
        if (!seen[abs(var)]) {
            seen[abs(var)] = True 

            if (trail[var_to_trail[abs(var)]].decisionLevel == curLevel) {
                levelCount = levelCount + 1
            }
        }
    }

    for (i = trail.size down to 0) {
        TrailEntry entry = trail[i]
        idx = abs(entry.lit)

        if (seen[idx]) {
            int[] reason = clauses[trail[i].reasonIdx]

            for each (int var in reason) {
                if (!seen[abs(var)]) {
                    seen[abs(var)] = True 

                    int trailIdx = var_to_trail[abs(var)]

                    if (trail[trailIdx].decisionLevel == curLevel) {
                        levelCount = levelCount + 1
                    }
                }
            }

            seen[idx] = False
            levelCount = levelCount - 1
        }

        if (levelCount == 1) {
            break
        }
    }

    int[] learned
    for (i = 1 up to numVars) {
        if (seen[i] == True) {
            int lit = trail[var_to_trail[i]].lit
            append(learned, -lit)
            seen[i] = False
        }
    }
}
\end{lstlisting}
\noindent
Here's the rundown of what's going on. In the first block, we're marking which 
variables we've seen and how many of them are from the current decision level from 
the initial conflicting clause. Next, we start traversing our trail from the most recent 
trail entry, and we move towards the front of the trail. If we've seen the variable from 
the current trail entry, we grab its reason clause, and for each new unseen variable, 
we mark it as seen and increment the current decision level count. Next, we ``remove''
the variable we just processed by unmarking it as seen and decrementing \textit{levelCount}.

Once we have only one literal from the current level, we construct a new learned clause 
by iterating over the entire \textit{seen} array and we add the negative literal to the 
learned clause because the learned clause is supposed to learn the opposite of what we've 
done so we don't repeat it. Then we unmark everything we have seen so that way we can reuse 
the same array the next time there's a conflict instead of allocating a new one each time. 

The results for this are very good, but I won't tell you what they are, because there's 
one more thing I want to do. You see, we iterate over the entire \textit{seen} array,
which is as large as the number of variables in our equation. However, it's very unlikely 
that we'll see many variables because first UIPs shrink how large the learned clause is, 
so there's less to be seen in the first place, but also because the more we learn, generally,
the learned clause gets smaller because each new learned clause adds more constrictions on 
which assignments are valid, and as we constrict which ones are valid and which ones aren't,
we inch towards learning the empty clause signaling that nothing is valid. This is just a 
long winded way to say that the iterating over the entire seen array each time, while O(n),
could be better.

But isn't that microoptimizing?
\footnote{Microoptimizing is like cutting your hair to help you reach a weight loss goal, 
instead of fixing your running form so you'll be able to have better cardio.}
After all, our two test sets have 20 and 50 variables 
respectively, which isn't much in the grand scheme of things. And you're right, our test 
sets are pretty small. However, the goal is to build a SAT solver that would scale well 
and most SAT solvers can handle problems with millions of variables pretty well, so I'd 
like this one to too. Making it scale well isn't the same as microoptimizing. Microoptimizing 
would be like saying \textit{++i} over \textit{i++}.
\footnote{Pre vs post incrementation in loops generally does not affect runtime and whether 
or not it actually matters is heavily debated.}

In order to do this, we're going to keep track of the indices of which variables we've
seen in an array and then iterate over that array. We will also need to keep track of 
which indices are in which position so we can swap and pop for efficient removal of 
elements.
\footnote{I love swap and pop if only for the name. I encourage everyone to do swap and 
pops so we can say ``swap and pop'' more than we currently say ``swap and pop.''}
Our object now looks like this:

\begin{lstlisting}
    public interface:
        SATSolver(string input)
        bool solveCNF()

    private interface:
        ...
        int[][] var_to_clause
        (int, int)[] clause_to_var

        bool[] seen
        int[] seenIdx
        int[] seenPos
\end{lstlisting}

\noindent
Here a subtle detail about these data structures. \textit{seen} and \textit{seenPos}
are a fixed size of \textit{numVars} because they require a full mapping from variable
to value at all times, but \textit{seenIdx} is dynamically sized because the number of 
variables that we see each clause resolution changes. So if you are implementing this 
yourself, make sure that \textit{seen} and \textit{seenPos} are allocated enough space 
upon initialization. Personally, I also reserve \textit{numVars} amount of space in 
\textit{seenIdx} so that it never has to reallocate space while solving, but it's not 
strictly necessary. 

This is how these data structures change the algorithm:
\begin{lstlisting}
if (contradiction has a value) {
    int[] conflict = contradiction.value
    int levelCount = 0
    int curLevel = decisionLevel

    for each (int var in conflict) {
        if (!seen[abs(var)]) {
            seen[abs(var)] = True 
            append(seenIdx, abs(var))
            seenPos[abs(var)] = seenIdx.size

            if (trail[var_to_trail[abs(var)]].decisionLevel == curLevel) {
                levelCount = levelCount + 1
            }
        }
    }

    for (i = trail.size down to 0) {
        TrailEntry entry = trail[i]
        idx = abs(entry.lit)

        if (seen[idx]) {
            int[] reason = clauses[trail[i].reasonIdx]

            for each (int var in reason) {
                if (!seen[abs(var)]) {
                    seen[abs(var)] = True 
                    append(seenIdx, abs(var))
                    seenPos[abs(var)] = seenIdx.size

                    int trailIdx = var_to_trail[abs(var)]

                    if (trail[trailIdx].decisionLevel == curLevel) {
                        levelCount = levelCount + 1
                    }
                }
            }

            seen[idx] = False

            //get the position of idx in seenIdx
            int pos = seenPos[idx]

            //get the last element in seenIdx
            int last = seenIdx.last

            //overwrite the element at the position pos with the value 
            //  of the last element in seenIdx
            seenIdx[pos] = last 

            //mark that the last element in seenIdx is now at position pos
            seenPos[last] = pos

            //remove the last element in seenIdx because it's also at position pos
            pop(seenIdx)

            levelCount = levelCount - 1
        }

        if (levelCount == 1) {
            break
        }
    }

    int[] learned
    while (seenIdx is not empty) {
        int lit_idx = seenIdx.last

        int lit = trail[var_to_trail[lit_idx]].lit
        append(learned, -lit)

        seen[lit_idx] = false
        pop(seenIdx)
        //do nothing with seenPos because we don't have to
    }
}
\end{lstlisting}

\subsection{Improvements - Results}

I'll let the results speak for themselves:
\noindent
uf20.91: \textbf{264ms}

\noindent
UUF50.218.1000: \textbf{2385ms}

\begin{FVerbatim}
On all 1000 equations in uf20-91:

Improved CDCL:             284ms (~0.28 seconds).
First UIPs:                406ms (~0.41 seconds).
Backjumping:               380ms (~0.38 seconds).
Clause Learning:           425ms (~0.43 seconds).
Iterative DPLL (AbP):      328ms (~0.33 seconds).
Iterative DPLL (PbA):      324ms (~0.32 seconds).
Watched Literals:          349ms (~0.35 seconds).
Pure Literal Elimination:  1548ms (~1.55 seconds).
Unit Propagation:          868ms (~0.87 seconds).
Recursive Backtracking:    1,540,613ms (~25.68 minutes).
Brute Force:               1,770,851ms (~29.51 minutes).
Short-Circuited (Basic):   49,874ms (~50 seconds).
Parallel:                  10,589ms (~10.5 seconds).

On all 1000 equations in UUF50.218.1000:

Improved CDCL:            2,385ms (~2.39 seconds).
First UIPs:               14,012ms (~14.01 seconds).
Backjumping:              11,291ms (~11.29 seconds).
Clause Learning:          9,652ms (~9.65 seconds).
Iterative DPLL (AbP):     6,369ms (~6.37 seconds).
Iterative DPLL (PbA):     6,365ms (~6.37 seconds).
Watched Literals:         7,130ms (~7.13 seconds).
Pure Literal Elimination: 231,107ms (~3.85 minutes).
Unit Propagation:         140,041ms (~2.34 minutes).
\end{FVerbatim}

\section{VSIDS}

Variable State Independent Decaying Sum (VSIDS) is a type of variable assignment 
heuristic, basically a loosely defined rule, that is used in modern CDCL solvers.
It is \textit{the} best variable assignment/branching heuristic with many variations.
In this chapter, we follow this paper
\begin{center}
    \url{https://arxiv.org/pdf/1506.08905}
\end{center}
which gives a much more in depth and theoretical explanation of how and why they work.
Four our purposes, we'll understand at a higher level what VSIDS are as well as how 
and why they work. 

Currently, our variable assignment choosing heuristic has been ``find first Unassigned
variable, and assign it True.'' With VSIDS, this becomes ``find the most active Unassigned
variable, and assign it True.'' The first step in VSIDS is to add in an array of
\textbf{activity scores} which rank ``how active'' a variable is. There are two schemes 
for ranking activity score, which is ranking each literal (positive and negative variables),
or just the variables (absolute value of the literal). For our purposes, we will rank just 
the variables. 

The next step is after clause resolution, add to each variable's score if it was involved 
in clause resolution. Again, there are multiple ways to do this. Some implementations choose 
only the variables in the learned clause, some implementations choose every variable involved 
in the resolution algorithm, even if they are resolved away, and there are more implementations 
than just those two. For our purposes, we will choose bumping the scores of only the variables 
in the learned clause.

The last step, which doesn't happen every iteration, is to multiply each score by a decay factor 
\footnote{Not to be confused with Deadpool's ``dying factor."}
after a certain number of conflicts. The idea behind this is to even the playing field for 
each variable. If a variable accumulates a giant activity score, obviously it will be the first 
pick if it's valid. However, activity scores don't represent how recent the activity score bump 
occurred, so if that variable stopped being active, it would still be the first pick despite 
not actually being active relative to the learned clauses. This decay lets the variables that 
are starting to become more active have a better chance at competing with the potentially 
inactive, but score-dominating variables. As it is a decaying factor, this value is chosen 
to be greater than 0, but less than 1, so that the scores shrink to some percentage.

\subsection{cVSIDS}

We're going to start out with implementing the most basic form of VSIDS and work our way 
up to some better versions. We'll start with the ``textbook'' form, which was outlined in 
the section's introduction. This form of VSIDS is called \textbf{cVSIDS} by the paper we're 
using. We first add in our \textit{activity} array into our solver 
object with a counter for the number of conflicts.
\begin{lstlisting}
    public interface:
        SATSolver(string input)
        bool solveCNF()

    private interface:
        ...
        double[] activity
        int numConflicts
        ...
\end{lstlisting}
Then our solve code becomes as follows:
\begin{lstlisting}
bool solve() {
    //setup
    ...

    while (True) {
        if (trailHead == trail.size) {
            double maxScore = -1.0
            optional int firstUnassigned

            for (i = 1 up to numVars) {
                if (assignment[i] == Unassigned) {
                    if (activity[i] > maxScore) {
                        firstUnassigned = i
                        maxScore = activity[i]
                    }
                }
            }  
            
            //the rest of the assignment code is the same
            ...
        }

        optional int[] contradiction = propagate(assignment)

        if (contradiction has a value) {
            numConflicts = numConflicts + 1

            //resolve conflict using UIPs
            ...
            int[] learned
            while (seenIdx is not empty) {
                int lit_idx = seenIdx.last
                activity[lit_idx] = activity[lit_idx] + 1
                ...
            }

            //backjump
            ...

            //apply decay factor
            if (numConflicts == X) {
                for (i = 0 up to activity.size) {
                    activity[i] *= Y
                }
                numConflicts = 0
            }

            continue
        }
    }
}
\end{lstlisting}
The pseudocode is very straightforward, so I won't bother to explain it much. Just 
so you know, if you choose to initialize \textit{maxScore} with a value of 0.0 in the 
assignment code, use \verb|>=| instead of \verb|>| because the activity scores start 
out at 0.0, so not using \verb|>=| will mess with unsatisfiable problems. In our 
code to apply the decay factor, I use ``X'' and ``Y'' to represent the number of 
conflicts and the decay factor respectively. The reason why these aren't set values,
but rather placeholders is because there aren't two numbers that perform the best for
\textit{all} equations, so they should be adjusted programmatically or manually set 
experimentally. However, I just decided to choose two random numbers, so I used 
\textbf{X = 10} and \textbf{Y = 0.8}. But results for different X and Y will be 
shown as well.

The results were as follows:
\noindent
uf20.91: \textbf{279ms}

\noindent
UUF50.218.1000: \textbf{1841ms}

We didn't get a big speedup for satisfiable problems, but we 
did shave off about half a second for the unsatisfiable problems, 
which I think is pretty darn good.

\begin{FVerbatim}
On all 1000 equations in uf20-91:

cVSIDS (X = 1, Y = 0.5):    260ms (~0.26 seconds).
cVSIDS (X = 1, Y = 0.8):    266ms (~0.27 seconds).
cVSIDS (X = 10, Y = 0.5):   270ms (~0.27 seconds).
cVSIDS (X = 10, Y = 0.8):   279ms (~0.29 seconds).
Improved CDCL:              284ms (~0.28 seconds).

On all 1000 equations in UUF50.218.1000:

cVSIDS (X = 1, Y = 0.5):    1,680ms (~1.68 seconds).
cVSIDS (X = 1, Y = 0.8):    1,705ms (~1.71 seconds).
cVSIDS (X = 10, Y = 0.5):   1,782ms (~1.78 seconds).
cVSIDS (X = 10, Y = 0.8):   1,841ms (~1.84 seconds).
Improved CDCL:              2,385ms (~2.39 seconds).
\end{FVerbatim}
\noindent
In the results, since we will be interating over many versions of VSIDS, 
I will leave out the older results and only compare against the previous 
versions of our solver. In the results, it appears to be that lower 
X and Y values lead to faster solving times, however, I encourage you to 
choose your own X and Y values and see what happens for yourself.

\subsection{mVSIDS}

\textbf{mVSIDS} is the second form of VSIDS introduced by the paper from 
the section introduction. In this type of VSIDS, we bump the values of 
every variable we resolve on during clause resolution as well as the values 
of the variables in the learned clause. The multiplicative decay step is 
applied the same as during cVSIDS.

The algorithmic change is basically trivial:
\begin{lstlisting}
bool solve() {
    //setup
    ..

    while (True) {
        //assigning
        ...

        optional int[] contradiction = propagate(assignment)

        if (contradiction has a value) {
            numConflicts = numConflicts + 1

            //set up UIPs from conflict clause
            ...

            for (i = trail.size down to 0) {
                TrailEntry entry = trail[i]
                idx = abs(entry.lit)

                if (seen[idx]) {
                    //grab reason and update seen and levelCount
                    ...

                    activity[idx] += 1
                    //unmark seen and update levelCount
                    ...
                }

                if (levelCount == 1) {
                    break
                }
            }

            //create learned clause and update activity scores
            ...

            //backjump
            ...

            //apply decay factor
            ...

            continue
        }
    }
}
\end{lstlisting}
\noindent 
Here are the results using \textbf{X = 10} and \textbf{Y = 0.8}:

\begin{FVerbatim}
On all 1000 equations in uf20-91:

mVSIDS (X = 1, Y = 0.5):    279ms (~0.28 seconds).
mVSIDS (X = 1, Y = 0.8):    277ms (~0.28 seconds).
mVSIDS (X = 10, Y = 0.5):   269ms (~0.27 seconds).
mVSIDS (X = 10, Y = 0.8):   284ms (~0.28 seconds).

cVSIDS (X = 1, Y = 0.5):    260ms (~0.26 seconds).
cVSIDS (X = 1, Y = 0.8):    266ms (~0.27 seconds).
cVSIDS (X = 10, Y = 0.5):   270ms (~0.27 seconds).
cVSIDS (X = 10, Y = 0.8):   279ms (~0.29 seconds).

Improved CDCL:              284ms (~0.28 seconds).

On all 1000 equations in UUF50.218.1000:

mVSIDS (X = 1, Y = 0.5):    1,532ms (~1.53 seconds).
mVSIDS (X = 1, Y = 0.8):    1,526ms (~1.53 seconds).
mVSIDS (X = 10, Y = 0.5):   1,420ms (~1.42 seconds).
mVSIDS (X = 10, Y = 0.8):   1,467ms (~1.47 seconds).

cVSIDS (X = 1, Y = 0.5):    1,680ms (~1.68 seconds).
cVSIDS (X = 1, Y = 0.8):    1,705ms (~1.71 seconds).
cVSIDS (X = 10, Y = 0.5):   1,782ms (~1.78 seconds).
cVSIDS (X = 10, Y = 0.8):   1,841ms (~1.84 seconds).

Improved CDCL:              2,385ms (~2.39 seconds).
\end{FVerbatim}
\noindent
So even with the worst mVSIDS time, it beats the best cVSIDS time
(per the benchmark, results may vary, these tests are very small and are not 
very representative of real-world SAT solver usage).

\subsection{EMA/Normalized VSIDS}

EMA stands for ``Exponential Moving Average'' and it's what allows us to 
mathematically decide which variables' activity scores should be decayed 
and by how much. If we go back to the earlier example of some variable gaining 
a really high activity score and then becoming inactive, but still dominating 
over the more recently active variables, this is what EMA helps with. EMA will 
decay high scoring, but actually inactive variables' scores while mostly 
preserving the actually active variables' scores. 

The way it does it is with this nifty formula:
$$s_n = (1 - f)\delta_n + f * s_{n - 1}$$

$s_n$ is the \textbf{current score}, $f$ is our \textbf{decay factor},
$\delta_n \in \{0, 1\}$ where a value of $0$ corresponds to a variable 
not being bumped, and a value of $1$ corresponds to a variable being bumped.
Let's walk through an example:

\begin{center}
    Suppose $f = 0.95$, and $s_0 = 0.0$. In the first conflict ($n = 1$),
    a variable was involved. Its score becomes:
    $$s_1 = (1 - 0.95)(1) + 0.95(0.0) = 0.05$$
    It's not much, but it's much more than 0, so this variable is considered 
    high ranking. Now suppose in the next conflict ($n = 2$), the variable was 
    not involved. Its score now becomes:
    $$s_2 = (1 - 0.95)(0) + 0.95(0.05) = 0.0475$$
    This is less than $0.05$, so now it's low ranking, even if the difference is small.
    Now suppose it's involved in the third conflict ($n = 3$):
    $$s_3 = (1 - 0.95)(1) + 0.95(0.0475) = 0.095125$$
    This is much larger than what it was before.
\end{center}

Although you probably have guessed already, we've been implementing VSIDS 
in a straightforward naive way by following the described process down to the letter. 
For sake of simplicity, for Normalized VSIDS, we'll also implement it naively.

Implementing normalized VSIDS isn't too different to how we're already implementing it,
except now, we update every score upon each conflict iteration. We'll need an additional 
data structure to keep track of the $\delta$ values. Since we're appending normalization 
on top of mVSIDS, we still update a $\delta$ value only when we resolve a variable or it 
appears in the learned clause. This means that we now delay the updating of our activity 
scores until after we construct the learned clause:
\begin{lstlisting}
bool solve() {
    //setup
    ..

    while (True) {
        //assigning
        ...

        optional int[] contradiction = propagate(assignment)

        if (contradiction has a value) {
            numConflicts = numConflicts + 1

            //set up UIPs from conflict clause
            ...

            for (i = trail.size down to 0) {
                TrailEntry entry = trail[i]
                idx = abs(entry.lit)

                if (seen[idx]) {
                    //grab reason and update seen and levelCount
                    ...

                    delta[idx] = True
                    //unmark seen and update levelCount
                    ...
                }

                if (levelCount == 1) {
                    break
                }
            }

            //create learned clause and update activity scores
            int[] learned
            while (seenIdx is not empty) {
                int lit_idx = seenIdx.last
                delta[lit_idx] = True
                ...
            }

            //backjump
            ...

            //apply decay factor
            for (i = 1 up to numVars) {
                if (delta[i] == True) {
                    activity[i] = (1 - Y) + Y * activity[i]
                } else {
                    activity[i] = Y * activity[i]
                }
                delta[i] = False
            }

            continue
        }
    }
}
\end{lstlisting}
\noindent
Notice how that there is no X value anymore because normalization happens upon 
every iteration instead of after a set number. Thankfully, our fastest times were 
when we updated activity scores after every iteration anyways, so we probably weren't 
missing some ``fastest'' time under a different X value.

Here are the results:
\begin{FVerbatim}
On all 1000 equations in uf20-91:

EMA (X = 1, Y = 0.95):      266ms (~0.27 seconds).
EMA (X = 1, Y = 0.8):       287ms (~0.29 seconds).
EMA (X = 1, Y = 0.5):       280ms (~0.28 seconds).

mVSIDS (X = 1, Y = 0.5):    279ms (~0.28 seconds).
mVSIDS (X = 1, Y = 0.8):    277ms (~0.28 seconds).
mVSIDS (X = 10, Y = 0.5):   269ms (~0.27 seconds).
mVSIDS (X = 10, Y = 0.8):   284ms (~0.28 seconds).

cVSIDS (X = 1, Y = 0.5):    260ms (~0.26 seconds).
cVSIDS (X = 1, Y = 0.8):    266ms (~0.27 seconds).
cVSIDS (X = 10, Y = 0.5):   270ms (~0.27 seconds).
cVSIDS (X = 10, Y = 0.8):   279ms (~0.29 seconds).

Improved CDCL:              284ms (~0.28 seconds).

On all 1000 equations in UUF50.218.1000:

EMA (X = 1, Y = 0.95):      1,504ms (~1.50 seconds).
EMA (X = 1, Y = 0.8):       1,497ms (~1.50 seconds).
EMA (X = 1, Y = 0.5):       1,621ms (~1.62 seconds).

mVSIDS (X = 1, Y = 0.5):    1,532ms (~1.53 seconds).
mVSIDS (X = 1, Y = 0.8):    1,526ms (~1.53 seconds).
mVSIDS (X = 10, Y = 0.5):   1,420ms (~1.42 seconds).
mVSIDS (X = 10, Y = 0.8):   1,467ms (~1.47 seconds).

cVSIDS (X = 1, Y = 0.5):    1,680ms (~1.68 seconds).
cVSIDS (X = 1, Y = 0.8):    1,705ms (~1.71 seconds).
cVSIDS (X = 10, Y = 0.5):   1,782ms (~1.78 seconds).
cVSIDS (X = 10, Y = 0.8):   1,841ms (~1.84 seconds).

Improved CDCL:              2,385ms (~2.39 seconds).
\end{FVerbatim}
\noindent
Here we see that EMA/normalizing the VSIDS doesn't make a huge contribution,
but it still makes a decent one for some values of Y.

\subsection{AMA/adaptVSIDS}

Adaptive Moving Average (AMA) VSIDS is an extension to EMA VSIDS where rather than 
having a fixed decay factor, we choose between two different decay rates depending on 
the state of our learned clause. The big idea behind making the decay factor variable 
instead of fixed is that our solver's ``productivity" isn't fixed. The learned clause 
affects this productivity in how high the quality of the learned clause is. 

To figure out the quality of a learned clause, we have to look at its \textbf{LBD}
which stands for \textbf{literal blocks distance}. LBD is defined as the number of 
distinct decision levels in the learned clause. A learned clause with a low LBD is considered to be higher 
quality than that of one with a high LBD. The paper doesn't explicitly say why this is the case,
\footnote{I might be illiterate and may have just missed the reason.}
but we can make a few guesses as to why.Considering that a low LBD means that there aren't 
many decision levels involved in the contradiction, so that means that the contradiction 
is more or less ``confined" to a smaller area of the search space. This indicates that 
the conflict has a more localized relationship with the current set of decisions. 
Such clauses have a higher chance to become unit or nearly unit, under a smaller 
amount of changes to the current assignment, better contraining future search more effectively.
Contrast this to high LBD, which has multiple decision levels involved in the contradiction, 
meaning that the contradiction spans a broader portion of the current search space, making 
it less likely it'll have an immediate effect on constraining future search.

This is supported by the fact that when we learn a clause with a low LBD, we choose a 
higher decay factor, so that variable activity scores shrink less, meaning we remember 
them more in the present. Likewise, when we learn a clause with a high LBD, we choose 
a lower decay factor, so variable activity scores shrink faster, making older activity 
lose its influence faster relative to the newer active variables.

The paper introduces this variable called \textit{lbdema}, which is defined as the 
exponential moving average of the LBD. Like our EMA VSIDS, this is calculated 
upon every iteration, and it's also a recurrence relation where the previous value affects 
what the next value will be. The formula is similar to the previous EMA formula:
$${lbdema}_n = (1 - \alpha) * {LBD}_t + \alpha * {lbdema}_{n - 1}$$
$\alpha$ is our \textbf{smoothing factor}, ${LBD}_t$ is the LBD of the current learned clause,
and \\${lbdema}_{n - 1}$ is the previous lbdema. In our testing of AMA VSIDS, I used 
$\alpha = 0.1$. 

\subsection{VSIDS - Results}
The results for VSIDS vary wildly because there are so many different types beyond what 
we've explored, and there are so many different parameters that could affect the end 
results. For VSIDS, although there might look like there's a clear winner, we need to 
keep in mind that we're using two relatively small test sets, so drawing an actual
conclusion as to which VSIDS version is ``objectively'' better isn't credible. 
Although AMA VSIDS isn't the fastest from our benchmarks, it will be the version we 
continue with for the rest of our paper. As for code, given how easy it is to swap out 
each version of VSIDS with each other, the exact code used for each version won't be 
provided, as the total changes were small enough that following along shouldn't be overly
difficult.

Another thing to note is that at the moment, there was not much exploration for the 
time differences. In order to figure out why the speeds between each version was 
different, it would be helpful to analyze how many conflicts there were, as well as the
order of deciding variables. As of now, I am just looking at the results and just 
accepting them as they are, not because I lack curiosity or scientific rigor, but because 
this is the second night in a row where I barely got any sleep for no reason, and I 
am too tired to analyze this under a microscope.

\begin{FVerbatim}
On all 1000 equations in uf20-91:

AMA (Y = 0.75 or 0.99):     274ms (~0.27 seconds).

EMA (X = 1, Y = 0.95):      266ms (~0.27 seconds).
EMA (X = 1, Y = 0.8):       287ms (~0.29 seconds).
EMA (X = 1, Y = 0.5):       280ms (~0.28 seconds).

mVSIDS (X = 1, Y = 0.5):    279ms (~0.28 seconds).
mVSIDS (X = 1, Y = 0.8):    277ms (~0.28 seconds).
mVSIDS (X = 10, Y = 0.5):   269ms (~0.27 seconds).
mVSIDS (X = 10, Y = 0.8):   284ms (~0.28 seconds).

cVSIDS (X = 1, Y = 0.5):    260ms (~0.26 seconds).
cVSIDS (X = 1, Y = 0.8):    266ms (~0.27 seconds).
cVSIDS (X = 10, Y = 0.5):   270ms (~0.27 seconds).
cVSIDS (X = 10, Y = 0.8):   279ms (~0.29 seconds).

Improved CDCL:              284ms (~0.28 seconds).

On all 1000 equations in UUF50.218.1000:

AMA (Y = 0.75 or 0.99):     1,629ms (~1.63 seconds).

EMA (X = 1, Y = 0.95):      1,504ms (~1.50 seconds).
EMA (X = 1, Y = 0.8):       1,497ms (~1.50 seconds).
EMA (X = 1, Y = 0.5):       1,621ms (~1.62 seconds).

mVSIDS (X = 1, Y = 0.5):    1,532ms (~1.53 seconds).
mVSIDS (X = 1, Y = 0.8):    1,526ms (~1.53 seconds).
mVSIDS (X = 10, Y = 0.5):   1,420ms (~1.42 seconds).
mVSIDS (X = 10, Y = 0.8):   1,467ms (~1.47 seconds).

cVSIDS (X = 1, Y = 0.5):    1,680ms (~1.68 seconds).
cVSIDS (X = 1, Y = 0.8):    1,705ms (~1.71 seconds).
cVSIDS (X = 10, Y = 0.5):   1,782ms (~1.78 seconds).
cVSIDS (X = 10, Y = 0.8):   1,841ms (~1.84 seconds).

Improved CDCL:              2,385ms (~2.39 seconds).
\end{FVerbatim}
\noindent

\section{Clause Database Management}

Clause database management is another essential addition to a CDCL SAT solver. 
For larger and larger equations, and for equations with less valid assignents, 
we will tend to learn more clauses. While First UIPs can efficiently compute 
these learned clauses, with the added benefit of them being smaller than if we 
didn't use the first UIP condition, the number of learned clauses still has a 
memory impact. Although our test sets are of 3-SAT instances, so the initial 
clauses are small to begin with, and the learned clauses aren't as large as they 
could be, for larger SAT problems, the number of learned clauses can become exponential
in size. 

Although some people like to say that memory management isn't important because the 
hardware gets better and better, the day we have to stop worrying about the amount 
of memory being used is the day we have unlimited memory to work with. And until that 
day comes (which will be never), those people can shut up and go back to coding their 
2GB Chrome tab website. 
\footnote{Websites take up too much memory, much more than they need, 
and I will happily die on this hill.}

For this section, we will go back to using this paper:
\begin{center}
    \small{\url{https://www3.cs.stonybrook.edu/~cram/cse505/Fall20/Resources/cdcl.pdf}}
\end{center}
starting from section \textbf{4.5.4. Clause Deletion Strategies}.
The three (listed) methods described are:
\begin{itemize}
    \item \textit{n-order learning}: Only learning clauses with $n$ or fewer literals.
    \item \textit{m-size relevance-based learning}: Learning all clauses as long as they 
            currently imply assignments or are themselves a unit clause, only being
            only being discarded once they have $m$ or more Unassigned literals.
    \item \textit{k-bounded learning}: Clauses smaller than $k$ are kept permanently,
            but larger clauses are discarded once they stop being unit.
\end{itemize}
\noindent
The paper also mentions that \textit{k-bounded learning} and \textit{m-size relevance-based learning}
can be combined. You may wonder why we can't combine all three, however, notice that 
\textit{k-bounded learning} is an extension of \textit{n-order learning}, as they both keep 
clauses smaller than some threshold, but \textit{k-bounded learning} is more merciful to 
larger clauses.

The paper also describes a newer heuristic that is based on its \textit{activity} score. 
The heuristic is essentially ``VSIDS for clauses instead of variables.'' We will progress our 
way into implementing all of these algorithms, incrementally of course. Before we continue, it 
should be noted that modern SAT solvers don't have a universal clause deletion policy, and each 
one does its own thing, and once ours is mature enough, it will be no different.
\footnote{Does that make us all unique or all the same?}

\subsection{N-Order Learning}

\textit{n-order learning} is definitely the most aggressive clause learning/deletion policy of 
the bunch. I don't think I'd want to encounter it on the streets. It's very straightforward in 
idea: if a clause has more than $n$ literals, it won't get learned at all. This may raise some red 
flags, and I hope the one we're all thinking of is, how do we backjump? If we're not learning a clause, 
how do we know which decision level to backjump to? Well, just because we're not learning it doesn't 
mean we can't backjump to a higher decision level, it just means we won't add the new clause to our 
database. Does that put you at ease? Well it shouldn't because the reason why this was such a big 
red flag in the first place is that as of now, when we decide a new variable, we always start with 
True and rely on clause learning to force it False later. So when we backjump and start a new 
decision level, how do we start the new decision level as False if we didn't learn anything?

As of now, there is no way to change our course.
\footnote{This will be addressed in the next section.}
So what can we do? You're gonna find this funny (not really), we need to bring back our 
``tried second branch'' flag, or at least some form of it. Without it, and I cannot emphasize 
this enough, we will just be stuck in an infinite loop.

So how should we pick a value of $n$? There isn't a universal way to choose, so we will try 
the maximum clause length of the original CNF, and then try some numbers around it. Because 
our current testing sets are all instances of 3-SAT, it's not very fair, in my opinion, to 
test on different sizes of $n$ when the data is specifically tailored to be a fixed size. Because 
of this, we will go to our CNF database 
\footnote{\url{https://www.cs.ubc.ca/~hoos/SATLIB/benchm.html}}
and get a new testing set, ``Flat200-479'' which is a SAT encoding of a 3-color problem. There are 
600 variables and 2237 clauses, all satisfiable.
\footnote{As I was setting up the benchmarking test, my framework reported that the solver failed on 
this new test set. It turned out I gave it the wrong path, and after providing it the correct path, it 
passed, as we would expect. The shock of seeing the [FAIL] really threw me for a loop.}

Here are the baseline results (without n-order learning):
\begin{FVerbatim}
On all 1000 equations in uf20-91:

AMA:     274ms (~0.27 seconds).

On all 100 equations in Flat200-479:

AMA:     24,238ms (~24.24 seconds).

On all 1000 equations in UUF50.218.1000:

AMA:     1,629ms (~1.63 seconds).
\end{FVerbatim}
\noindent
We will now discuss our changes to our algorithm. Firstly, we need to compute the 
maximum clause length. As of now, we will say that the maximum clause length of the 
original equation will also be the maximum clause length of any learned clause, and then 
see how changing this number to lower or higher changes our runtime speed. To compute 
the maximum clause length, we will iterate through the original clauses and take the 
maximum size.

Next, we want to check how large the learned clause will be. We do this by checking 
the size of the \textit{seenIdx} array because it holds the variables that will 
become part of the learned clause. This avoids computing what the learned clause will 
be unnecessarily. For now, we will just backtrack one level in this case, then move onto 
backjumping.

We first need to remember to change our interface to use our old DecisionLevel struct:
\begin{lstlisting}
    public interface:

    private interface:
        DecisionLevel {
            int idx,
            bool triedOtherBranch
        }

        DecisionLevel[] trailStarts
        ...
\end{lstlisting}
\noindent
And of course we need to adjust our code to use the DecisionLevel struct instead of just an index.

The algorithmic changes are as follows:
\begin{lstlisting}
bool solve() {
    //setup
    ..

    while (True) {
        //assigning
        ...

        optional int[] contradiction = propagate(assignment)

        if (contradiction has a value) {
            numConflicts = numConflicts + 1

            //set up UIPs from conflict clause
            ...

            for (i = trail.size down to 0) {
                TrailEntry entry = trail[i]
                idx = abs(entry.lit)

                if (seen[idx]) {
                    //grab reason and update seen and levelCount
                    ...

                    delta[idx] = True
                    //unmark seen and update levelCount
                    ...
                }

                if (levelCount == 1) {
                    break
                }
            }

            if (seenIdx.size > maxSize) {
                while (trail.size > 0 && trail_starts.size > 0) {
                    DecisionLevel decision = trailStarts.last

                    while (trail.size > decision.idx + 1) {
                        int lit = trail.last.lit
                        assignment[abs(lit)] = Unassigned
                        var_to_trail[abs(lit)] = -1
                        erase(trail, trail.last)
                    }

                    if (decision.triedFalse == False) {
                        TrailEntry entry = trail.last

                        entry.lit = -entry.lit
                        assignment[abs(entry.lit)] = False
                        var_to_trail[abs(entry.lit)] = -1
                        decision.triedFalse = True

                        trailHead = decision.idx
                        break
                    } else {
                        int lit = trail.last.lit
                        assignment[abs(lit)] = Unassigned
                        var_to_trail[abs(lit)] = -1
                        erase(trail.last)
                        erase(trailStarts.last)
                    }
                    
                    if (trail.size == 0) {
                        return False
                    }

                    while (seenIdx.size > 0) {
                        int idx = seenIdx.last
                        seen[idx] = False
                        pop(seenIdx)
                    }

                    for (i = 0 up to numVars) {
                        delta[i] = False
                    }

                    continue
                }
            }

            //create learned clause and update activity scores
            int[] learned
            while (seenIdx is not empty) {
                int lit_idx = seenIdx.last
                delta[lit_idx] = True
                ...
            }

            //backjump
            ...

            //apply decay factor
            for (i = 1 up to numVars) {
                if (delta[i] == True) {
                    activity[i] = (1 - Y) + Y * activity[i]
                } else {
                    activity[i] = Y * activity[i]
                }
                delta[i] = False
            }

            continue
        }
    }
}
\end{lstlisting}
\noindent
So it's pretty straightforward,
\footnote{I struggled with debugging for half an hour because I forgot to clear out
\textit{delta} and \textit{seenIdx/seen}. I can see the headlines now: 
``Functional programmers are right about this ONE thing!!!''}
we separate our code into two branches. In the first branch where we aren't going to 
learn the clause, we just backtrack as we used to in our clause learning section, clear 
out everything related to adjusting the VSIDS and then move on, and otherwise we move on 
business as usual.

\begin{FVerbatim}
On all 1000 equations in uf20-91:

N-order:   283ms (~0.28 seconds).
AMA:       274ms (~0.27 seconds).

On all 100 equations in Flat200-479:

N-order:   85,749ms (~1.43 minutes).
AMA:       24,238ms (~24.24 seconds).

On all 1000 equations in UUF50.218.1000:

N-order:   2,313ms (~2.31 seconds).
AMA:       1,629ms (~1.63 seconds).
\end{FVerbatim}
\noindent
As you can see, doing n-order learning was terrible for our runtime. Unfortunately 
this isn't an implementation issue in our n-order learning algorithm, unless you 
consider that implementing n-order learning is an implementation issue to begin with.
You can probably guess why this is (we removed a ton of unit propagation!), and this 
is really important because it goes to show that even small initial clauses can cause 
much larger learned clauses. When I doubled $n$, the runtime was nearly the same as 
AMA, but a tad slower. 

\subsection{M-Size Relevance-Based Learning}

The kind of clause management heuristic we're going to try is m-size relevance-based 
learning, where we learn every single clause, but only delete them when they stop 
implying assignments and have more than $m$ Unassigned literals. To set this up, 
we will restart from our VSIDS code since it's not like there's anything to build off of 
from n-order learning. The paper we use uses the phrasing ``as soon as the number of
unassigned literals is greater than an integer $m$'' which sounds like we have to keep 
track of the number of unassigned literals per clause. We can do this by directly tying 
the number of Unassigned literals to the clause in a struct, which we will do. 

There is also the option of scanning the clauses in their entirety and using that information 
to choose which ones to keep and which ones to discard. Then, instead of doing this every 
conflict or whatever, we clean up the database after a certain number of conflicts.

Although adding in this struct to keep track of the number of Unassigned literals seems 
really simple, you have to consider that our watched literals implementation isn't a great 
support structure for it, since we'd need to add in a mapping from variables to clauses,
then update every single clause containing a variable, and blah blah blah, it's a lot.
There's also just the fact that honestly, I'm going to be debugging this for hours, the 
end of the quarter is draining me and I just want to graduate, and so in other words,
I'm making the executive decision to not ``waste'' more time than I have to.

So since our clause lengths are pretty small still, we'll say that once the clause 
is no longer unit, we'll get rid of it, as long as it's not part of forcing an assignment.
The implementation will be extensible, ie, it won't ``hardcode'' that the clause must be 
non-unit in order to be deleted, it will still manually count the number of Unassigned 
literals and then compare against a threshold to decide.

Here's what we need to do to get this to work: every time there's a certain number 
of conflicts, we need to scan each clause if they meet our conditions. Then for 
each ones that do, we will mark them to be deleted. Then we just delete them! So simple!
So easy! It's actually more complicated than that because of our watched literals unfortunately.
We need to remove the watch lists of each clause we're deleting, update the watch lists of all 
variables that were watching the clause we deleted, and update the watch lists of all the clauses 
that get shifted in memory, which of course, is a lot easier said than done. So here's the game 
plan: I'm gonna \textit{try} to get some sleep (it is 1:41AM), and when I wake up, maybe I'll 
have the mental fortitude to properly explain this.
\footnote{My senioritis got so bad, it evolved into full blown mental illness.}

Here's what we'll need first: we're going to implement the helper functions to help update 
our watch lists so our watched literals will work properly. So first we'll add them to our 
object:

\begin{lstlisting}
    public interface:

    private interface:
        ...
        void removeFromWatchList(int[] watchList, int idx)
        void removeClauseFromWatchLists(int clauseIdx)
        void replaceWatchIndex(int oldIdx, int newIdx)
        ...
\end{lstlisting}
\noindent
We will implement them in order.

\begin{lstlisting}
void removeFromWatchList(int[] watchList, int idx) {
    for (i = 0 up to watchList.size) {
        if (watchList[i] == idx) {
            swap(watchList[i], watchList.last)
            pop(watchList)
            break
        }
    }
}
\end{lstlisting}
This algorithm is very straightforward. Given an array of clause indices that 
we're watching, we want to remove the passed in \textit{clause index}, not 
the \textit{element} at the passed in index. We use swap-and-pop to avoid 
shifting elements, which would be a bookkeeping nightmare.

\begin{lstlisting}
void removeClauseFromWatchLists(int clauseIdx) {
    (int firstWatch, int secondWatch) = clause_to_var[clauseIdx]

    int lit1 = clauses[clauseIdx][firstWatch]
    int lit2 = clauses[clauseIdx][secondWatch]

    int varIdx1 = var_to_widx(lit1)
    int varIdx2 = var_to_widx(lit2)

    removeFromWatchList(var_to_clause[varIdx1], clauseIdx)

    if (varIdx1 != varIdx2) {
       removeFromWatchList(var_to_clause[varIdx2], clauseIdx)
    }
}
\end{lstlisting}
\noindent
This helper works by first grabbing the indices of the two literals watching 
the clause at the passed in index. We then turn them into the actual literals 
by indexing the clause. Then, we get the indices of the literals to where they 
map in \textit{var\_to\_clause}, so that way we can grab the literals' watchlists.
Finally, we remove the clause index from those literals' watchlists. This is a 
little bit confusing because of all of the indexing and grabbing the corresponding 
elements we're doing. 

\begin{lstlisting}
void replaceWatchIndex(int oldIdx, int newIdx) {
    (int firstWatch, int secondWatch) = clause_to_var(newIdx)

    int lit1 = clauses[newIdx][firstWatch]
    int lit2 = clauses[newIdx][secondWatch]

    int varIdx1 = var_to_widx(lit1)
    int varIdx2 = var_to_widx(lit2)

    int[] firstWatchList = var_to_clause[var1]
    for (i = 0 up to firstWatchList.size) {
        if (firstWatchList[i] == oldIdx) {
            firstWatchList[i] = newIdx
        }
    }

    if (varIdx1 != varIdx2) {
        int[] secondWatchList = var_to_clause[var2]
        for (i = 0 up to secondWatchList.size) {
            if (secondWatchList[i] == oldIdx) {
                secondWatchList[i] = newIdx
            }
        }
    }
}
\end{lstlisting}
\noindent
This one is pretty self explanatory. We do our indexing magic to get the watchlists 
for the clause's corresponding watch literals, and we replace an old index with 
a new index. Again, not an element \textit{at} an old index, an element \textit{equaling}
an old index.

Here is how they are used in our m-size learning algorithm:
\begin{lstlisting}
bool solveCNF() {
    //setup
    ...
    int numConflicts = 0

    while(true) {
        if (numConflicts > threshold) {
            bool[] remove

            for (i = numClauses up to clauses.size) {
                int[] clause = clauses[i]
                
                int numUnassigned = 0
                for (j = 0 up to clause.size) {
                    if (vars[abs(clause[j])] == Unassigned) {
                        numUnassigned = numUnassigned + 1
                    }
                }

                if (numUnassigned > m) {
                    remove[i - numClauses] = True
                }
            }

            for (i = 0 up to trail.size) {
                if (trail[i].reasonIdx is not optional) {
                    int reasonIdx = trail[i].reasonIdx.value

                    if (reasonIdx >= numClauses) {
                        remove[reasonIdx - numClauses] = False
                    }
                }
            }

            for (i = clauses.size down to numClauses) {
                if (remove[i - numClauses] == False) {
                    continue
                }

                int lastIdx = clauses.size - 1

                if (i == lastIdx) {
                    removeClauseFromWatchLists(i)
                    pop(clauses)
                    pop(clause_to_var)
                    continue
                }

                removeClauseFromWatchLists(i)
                swap(clauses[i], clauses[last])
                swap(clause_to_var[i], clause_to_var[last])

                replaceWatchIndex(last, idx)

                for (i = 0 up to trail.size) {
                    if (trail[i].reasonIdx is not optional) {
                        if (trail[i].reasonIdx.value == last) {
                            trail[i].reasonIdx.value = idx
                        }
                    }
                }

                pop(clauses)
                pop(clause_to_var)
            }

            numConflicts = 0
        }

        //rest of solving loop
        ...
    }
} 
\end{lstlisting}
\noindent
This is what's going on: first, we check that a certain number of conflicts have happened. 
This number is fixed. If there have been that many conflicts, we scan every clause that is 
\textbf{not} a part of our original CNF equation. If it has more than $m$ Unassigned literals, 
we mark it for deletion. In our code, our \textit{remove} array uses an indexing where we subtract 
the original number of clauses from the current clause indexing variable. This keeps the array 
smaller in our implementation. Next, we iterate through our trail, and we check if any of the 
clauses we've marked to be deleted are the reason why something was propagated. If a clause we 
want to delete is a reason clause, we can't delete it because it's still being used in our 
assignments. 

Next, we iterate from the back of the clause database up to the front, excluding the 
original clauses. If we're not supposed to remove the current clause, we move on (duh).
Otherwise, we grab the index of the last clause. If the last clause in the database is 
marked to be deleted, we can easily pop it off the database. Otherwise, we need to do 
our swap and pop nonsense. We remove the clause from its watchlists, swap the current 
clause with the last clause in the database, then we replace the watch indices. 

After that, we traverse our trail again and check if we've swapped any literal's 
reason clause index. If we have, then we need to update it. After that, we can 
remove the clause and its watched literals from the database. After this is all over,
we reset the counter for the number of conflicts and continue with the rest of the 
solving loop. 

That's a lot of work, so how does it hold up against our other implementations?
\begin{FVerbatim}
On all 1000 equations in uf20-91:

M-size:    293ms (~0.29 seconds).
N-order:   283ms (~0.28 seconds).
AMA:       274ms (~0.27 seconds).

On all 100 equations in Flat200-479:

M-size:    25,681ms (~25.68 seconds).
N-order:   85,749ms (~1.43 minutes).
AMA:       24,238ms (~24.24 seconds).

On all 1000 equations in UUF50.218.1000:

M-size:    1,937ms (~1.94 seconds).
N-order:   2,313ms (~2.31 seconds).
AMA:       1,629ms (~1.63 seconds).
\end{FVerbatim}
\noindent
We see that m-size relevance-based learning is slightly slower than AMA VSIDS, and much faster 
than n-order learning (except for uf20-91). There could be some reasons for this, such as doing 
clause deletion too frequently. Rather than doing it a fixed number of times, we could also 
perform it whenever the database grows by a certain amount, say 10\%. If we do this, our results 
speed up quite a bit on the larger equations (FLAT200.470):

\begin{FVerbatim}
On all 1000 equations in uf20-91:

M-size (10\%):  320ms (~0.32 seconds).
M-size:         293ms (~0.29 seconds).
N-order:        283ms (~0.28 seconds).
AMA:            274ms (~0.27 seconds).

On all 100 equations in Flat200-479:

M-size (10\%):  17,716ms (~17.72 seconds).
M-size:         25,681ms (~25.68 seconds).
N-order:        85,749ms (~1.43 minutes).
AMA:            24,238ms (~24.24 seconds).

On all 1000 equations in UUF50.218.1000:

M-size (10\%)   1,768ms (~1.77 seconds).
M-size:         1,937ms (~1.94 seconds).
N-order:        2,313ms (~2.31 seconds).
AMA:            1,629ms (~1.63 seconds).
\end{FVerbatim}
\noindent
Although it doesn't speed up on the smaller equations, a ~7 second speedup on larger 
equations shows that this scales better on larger SAT problems, which is great!

\subsection{Combining M-Size and K-Bounded Learning}

Due to the similarities of m-size learning and k-bounded learning, I felt like it 
didn't make much sense to implement k-bounded learning by itself, then combining it 
with m-size learning afterwards. Combining these two is really easy since all we're 
doing is checking if clauses are:
\begin{enumerate}
    \item Greater than some size $k$
    \item Have at least $m$ Unassigned literals
\end{enumerate}
Since we already do the second one in our m-size loop, it's just adding another condition
that needs to be true in order to mark a clause to be removed.

On the other hand, what isn't easy is choosing a size $k$. Not all CNF equations are 
made equal. Some equations have small clauses (like our testing sets) and learn larger 
clauses, some have large clauses and learn smaller clauses, etc. Choosing a fixed $k$ for 
all clauses is obviously a really bad idea. Make it too large and we just learn everything, 
make it too small and we forget everything. We will bring back our ``maxSize'' counter 
from n-order based learning and use that. We will use a maxSize that is twice as large as 
the largest clause in the original CNF
\footnote{Because adding a second condition to one conditional is easy, and iterating through 
the clauses to find the size of the largest clause is also easy, we will skip pseudocode.}

\begin{FVerbatim}
On all 1000 equations in uf20-91:

MK-bound:       288ms (~0.29 seconds).
M-size (10\%):  320ms (~0.32 seconds).
M-size:         293ms (~0.29 seconds).
N-order:        283ms (~0.28 seconds).
AMA:            274ms (~0.27 seconds).

On all 100 equations in Flat200-479:

MK-bound:       19,082ms (~19.08 seconds).
M-size (10\%):  17,716ms (~17.72 seconds).
M-size:         25,681ms (~25.68 seconds).
N-order:        85,749ms (~1.43 minutes).
AMA:            24,238ms (~24.24 seconds).

On all 1000 equations in UUF50.218.1000:

MK-bound:       1,740ms (~1.74 seconds).
M-size (10\%)   1,768ms (~1.77 seconds).
M-size:         1,937ms (~1.94 seconds).
N-order:        2,313ms (~2.31 seconds).
AMA:            1,629ms (~1.63 seconds).
\end{FVerbatim}
\noindent
We see that using a maxSize that's twice as large as the size of the largest original clause
causes a pretty big slowdown on Flat200-479. This indicates that we're removing clauses 
too aggressively relative to the size of the clauses we're learning.

\subsection{``Clause VSIDS''}

Now we get to the more interesting part: ``Clause VSIDS.''
\footnote{This is NOT an official or technical name as far as I am aware}
The idea is that just like in VSIDS for variables, we measure a clause's activity
and quality (LBD), and use that to determine whether or not it should be deleted. 
We can still take into account things such as length, whether or not it's being used, 
and the other factors we've been considering in other versions as well, however, the 
activity and quality are the main ones.

The paper defines this heuristic as ``the decision whether a clause should be 
deleted is based not only on the number of literals but also on its \textit{activity} in 
contributing to conflict making and on the number of decisions taken since its creation.''

In order to implement this, we'll need three additional data structures:
\begin{lstlisting}
    public interface:

    private interface:
        ...
        double[] cActivity
        int[] creationConflict
        int[] lbds
        ...
\end{lstlisting}
\noindent
\textit{cActivity} measures a learned clause's activity. We will update the activity of 
every clause involved in a conflict, if it's a learned clause. Because activity is a metric 
we use to determine whether or not a clause should be deleted, and we never delete an original 
clause, we don't need to keep track of any data for any original clauses. 
\textit{creationConflict} keeps track of the conflict number that a clause was created on.
So if we learn a clause on conflict 50, it's creation conflict is 50. We use this to measure the 
age of the clause, where $age = totalConflicts - creationConflict$.
Lastly, \textit{lbds} keeps track of a clause's LBD upon learn time. Even though a clause's LBD 
changes as our decision trail changes, we want to keep the original LBD from when the clause was 
learned. This is because we want to know what decision levels led to the conflict in the first place, 
and recomputing it describes the how ``useful'' the learned clause is under the current assignment, 
not when it was actually created. Some SAT solvers do use dynamic LBD computations, but it's a 
different heuristic.

In order to keep track of these metrics for each learned clause, we need to update our 
VSIDS code, as well as some other parts of our conflict resolution algorithm. Luckily, 
these changes are very minimal:
\begin{lstlisting}
bool solveCNF() {
    //setup
    ...
    int totalConflicts = 0
    int curSize = numClauses

    double activityIncrease = 1.0
    double activityGrowth = 1.05
    double maxIncrease = 1e100
    double resetIncrease = 1e-100

    while (true) {
        if (activityIncrease > maxIncrease) {
            for each (clauseActivity in cActivity) {
                clauseActivity *= resetIncrease
            }
            activityIncrease *= resetIncrease
        }

        //Current clause deletion
        if (clauseDatabase.size grew by X%) {
            ...

            for (i = clauses.size down to numClauses) {
                if (remove[i - numClauses] == False) {
                    continue
                }

                int lastIdx = clauses.size - 1

                if (i == lastIdx) {
                    removeClauseFromWatchLists(i)
                    pop(clauses)
                    pop(clause_to_var)

                    pop(cActivity)
                    pop(creationConflict)
                    pop(lbds)

                    continue
                }

                removeClauseFromWatchLists(i)
                swap(clauses[i], clauses[last])
                swap(clause_to_var[i], clause_to_var[last])

                swap(cActivity[i], cActivity[last])
                swap(creationConflict[i], creationConflict[last])
                swap(lbds[i], lbds[last])

                ...

                pop(clauses)
                pop(clause_to_var)

                pop(cActivity)
                pop(creationConflict)
                pop(lbds)
            }

            curSize = clauses.size
        }

        //assignment
        ...

        if (contradiction occurred) {
            totalConflicts = totalConflicts + 1
            ...

            for (int trailIdx = trail.size down to 0) {
                ...
                int reasonIdx = trail[trailIdx].reasonIdx.value;
                if (reasonIdx >= numClauses) {
                    cActivity[reasonIdx - numClauses] += activityIncrease
                }
            }
            //rest of VSIDS and LBD calculation
            ...

            //backjump and learn 
            ... 

            append(cActivity, 0.0)
            append(creationConflict, totalConflicts)
            append(lbds, lbd)
            activityIncrease *= activityGrowth
            ...
        }
    }
}
\end{lstlisting}
\noindent
In our clause deletion code, we just added code to update our new data structures
to be consistent with the rest of the clause database data structures. And then 
in the VSIDS calculation, we added code to update a clause's activity if it's 
a learned clause, and at the end of the loop, we just add extra entries into each 
of the data structures so they grow with the database. 

We use an activity increase and an activity growth variable to simulate the effects 
of an activity decay and bump. The idea is that instead of bumping only the clauses 
that were involved in a conflict and decaying the rest, we bump only the clauses 
involved in a conflict in a way that is as if we decayed the others. By increasing 
the activity increase by some growth amount every conflict, when we add to a clause's 
activity, we add an ever increasing amount, so that way, more recent clauses get 
larger and larger activities and less recent clauses don't have such huge increases. 
The only caveat to this is that because conflicts happen all the time, we need to 
reset the activity increase once it grows to a stupidly large size. This can be thought 
of as our decay step.

The next step is to now use these metrics in our clause deletion code. We could do 
a bunch of experiments of how each metric impacts runtime and clause deletion, but 
frankly, I don't want to. Instead, here's what workflow we'll use:
\begin{enumerate}
    \item Find which clauses that can be deleted by LBD, age, and activity
    \item Find which clauses can be deleted via m+k-bounded learning
    \item Protect clauses with an LBD \verb|<=| 2
    \item Protect clauses that are currently reason clauses
    \item Delete the ``bad'' ones
\end{enumerate}
\noindent
To determine what a ``bad'' clause means, we need to think of what the goals of 
clause deletion is. For instance, do we want to delete a fixed number of clauses, 
or should it be more dynamic, as in, there isn't a quota to fulfill, instead we look at 
each clause and go ``you suck, get out'' or ``you're good, stay?'' This idea is pretty 
attractive, but there is something we should consider: what if there's not enough bad 
clauses relative to our database size, or what if there's too many bad clauses relative 
to our database size? Ie, what if we mark too little clauses for deletion and our database 
is still large, as well as marking too many clauses for deletion and then our database 
becomes super small again? To be honest, if there aren't enough ``bad'' clauses, to delete,
to me it seems like most of the database is really good, so I wouldn't mark \textit{extra}
clauses for deletion, however, capping how many clauses we delete does seem like a good idea.

Because we have a clear marker of when clause deletion should occur (when the database 
grows by 10\%), we will say that the number of clauses being deleted can't exceed this 
10\% threshold. So if our original database has 100 clauses and it grows to 110 clauses,
we can't delete more than 9 clauses. Supposing we delete 0, and the database grows to 
121, we can't delete more than 10 clauses, and so on. Another advantage of this is 
that the maximum number of clauses we can delete grows/shrinks with our threshold, so 
if in a later experiment we found that it's more beneficial to run the deletion algorithm 
every time the database grows by 20\%, we don't have to manually update the number of 
clauses we're allowed to delete as well.

In order to decide which clauses to delete, we want to decide a threshold for our 
metrics. I used the ones in the bottom 25th percentile for their respetive metric, 
so high LBD, high age, and low activity. We then select the clauses in those 
bottom metrics, and if they have at least 2 of these negative attributes, we mark 
them to be deleted. 

After that, we do our regular m+k bounded learning, then afterwards, we mark clauses 
that have low LBD (\verb|<=| 2), and clauses that are reason clauses as protected, 
so they can't be deleted even if they had at least two bad attributes.

The code is as follows:
\begin{lstlisting}
bool solveCNF() {
    //setup
    ... 

    while (true) {
        //Current clause deletion
        if (clauseDatabase.size grew by X%) {
            int numCanDelete = clauseDatabase.size - curSize

            int[] candidateIdx
            for (i = 0 up to clauseDatabase.size - numClauses) {
                candidateIdx[i] = i
            }

            int cutoff = candidateIdx.size * 0.25

            sortIncreasing(cActivity)

            double activityCutoff = cActivity[cutoff]

            sortIncreasing(conflictCreation)

            int ageCutoff = conflictCreation[cutoff]

            sortDecreasing(lbds)

            int lbdCutoff = lbds[cutoff]

            for (i = numClauses up to clauses.size) {
                bool lowActivity = cActivity[i] <= activityCutoff
                bool highAge = creationConflict[i] <= ageCutoff
                bool highlbd = lbds[i] >= lbdCutoff

                remove[i - numClauses] = lowActivity + highAge + highlbd >= 2
            }

            //rest of the database deletion
            ...
        }

        //rest of solving
        ...
    }

}
\end{lstlisting}

\begin{FVerbatim}
On all 1000 equations in uf20-91:

Clause VSIDS:   347ms (~0.35 seconds).
MK-bound:       288ms (~0.29 seconds).
M-size (10\%):  320ms (~0.32 seconds).
M-size:         293ms (~0.29 seconds).
N-order:        283ms (~0.28 seconds).
AMA:            274ms (~0.27 seconds).

On all 100 equations in Flat200-479:

Clause VSIDS:   19,181ms (~19.18 seconds).
MK-bound:       19,082ms (~19.08 seconds).
M-size (10\%):  17,716ms (~17.72 seconds).
M-size:         25,681ms (~25.68 seconds).
N-order:        85,749ms (~1.43 minutes).
AMA:            24,238ms (~24.24 seconds).

On all 1000 equations in UUF50.218.1000:

Clause VSIDS:   1,804ms (~1.80 seconds).
MK-bound:       1,740ms (~1.74 seconds).
M-size (10\%)   1,768ms (~1.77 seconds).
M-size:         1,937ms (~1.94 seconds).
N-order:        2,313ms (~2.31 seconds).
AMA:            1,629ms (~1.63 seconds).
\end{FVerbatim}
\noindent
To be honest, I'm a little bit disappointed in these results. I was really hopeful 
that this heuristic would have yielded more gains in speed, not runtime. Was my 
implementation perfect? No. Could that have contributed to the increased runtime?
Yes, but I have a feeling that it's also the testing sets being used. Again, we 
can't necessarily draw a conclusion on \textit{all} SAT problems based on this 
data, but we can draw the conclusion that the solve is worse for \textit{these}
SAT problems.

There were some things I tried to see what could be causing this. I tried adjusting 
how aggressive/loose the badness score had to be in order to be marked to be deleted. 
A lower badness score threshold correlated with a higher runtime. Then I tried removing 
the cap on deletion altogether, and that didn't impact the runtime very much. Then I tried 
adjusting the percentile to be the 20th, 25th, and 50th, and they all had very similar runtimes 
as well. I even tried messing with what we consider a bad clause, and that impacted 
the runtime, in a bad way. Lastly, I tried removing mk-bound learning, and that also didn't 
make much of an impact. I then removed the k-bounded learning afterwards, and that also didn't 
seem to help. No matter what I tried, it appears that this heuristic was hurting the runtime 
on all accounts. Perhaps something to be explored at a later date.

\subsection{Clause Database Management - Results}

Here are the results for all the kinds of clause database management compared to 
the last implementation without any kind of management:

\begin{FVerbatim}
On all 1000 equations in uf20-91:

Clause VSIDS:               347ms (~0.35 seconds).
MK-bound:                   288ms (~0.29 seconds).
M-size (10\%):              320ms (~0.32 seconds).
M-size:                     293ms (~0.29 seconds).
N-order:                    283ms (~0.28 seconds).

AMA:                        274ms (~0.27 seconds).
EMA (X = 1, Y = 0.95):      266ms (~0.27 seconds).

Improved CDCL:              284ms (~0.28 seconds).
First UIPs:                 406ms (~0.41 seconds).
Backjumping:                380ms (~0.38 seconds).
Clause Learning:            425ms (~0.43 seconds).
Iterative DPLL (AbP):       328ms (~0.33 seconds).
Iterative DPLL (PbA):       324ms (~0.32 seconds).
Watched Literals:           349ms (~0.35 seconds).
Pure Literal Elimination:   1548ms (~1.55 seconds).
Unit Propagation:           868ms (~0.87 seconds).
Recursive Backtracking:     1,540,613ms (~25.68 minutes).
Brute Force:                1,770,851ms (~29.51 minutes).
Short-Circuited (Basic):    49,874ms (~50 seconds).
Parallel:                   10,589ms (~10.5 seconds).

On all 100 equations in Flat200-479:

Clause VSIDS:   19,181ms (~19.18 seconds).
M-size (10\%):  17,716ms (~17.72 seconds).
AMA:            24,238ms (~24.24 seconds).

On all 1000 equations in UUF50.218.1000:

Clause VSIDS:               1,804ms (~1.80 seconds).
MK-bound:                   1,740ms (~1.74 seconds).
M-size (10\%)               1,768ms (~1.77 seconds).
M-size:                     1,937ms (~1.94 seconds).
N-order:                    2,313ms (~2.31 seconds).

AMA:                        1,629ms (~1.63 seconds).
EMA (X = 1, Y = 0.95):      1,504ms (~1.50 seconds).

Improved CDCL:              2,385ms (~2.39 seconds).
First UIPs:                 14,012ms (~14.01 seconds).
Backjumping:                11,291ms (~11.29 seconds).
Clause Learning:            9,652ms (~9.65 seconds).
Iterative DPLL (AbP):       6,369ms (~6.37 seconds).
Iterative DPLL (PbA):       6,365ms (~6.37 seconds).
Watched Literals:           7,130ms (~7.13 seconds).
Pure Literal Elimination:   231,107ms (~3.85 minutes).
Unit Propagation:           140,041ms (~2.34 minutes).
\end{FVerbatim}
\noindent

Because of how well M-size (10\%) performed, I will keep it in the overall results 
when compared with every other solver we've created so far:

\begin{FVerbatim}
On all 1000 equations in uf20-91:

Clause VSIDS:               347ms (~0.35 seconds).
M-size (10\%):              320ms (~0.32 seconds).

AMA (Y = 0.75 or 0.99):     274ms (~0.27 seconds).
EMA (X = 1, Y = 0.95):      266ms (~0.27 seconds).

Improved CDCL:              284ms (~0.28 seconds).
First UIPs:                 406ms (~0.41 seconds).
Backjumping:                380ms (~0.38 seconds).
Clause Learning:            425ms (~0.43 seconds).
Iterative DPLL (AbP):       328ms (~0.33 seconds).
Iterative DPLL (PbA):       324ms (~0.32 seconds).
Watched Literals:           349ms (~0.35 seconds).
Pure Literal Elimination:   1548ms (~1.55 seconds).
Unit Propagation:           868ms (~0.87 seconds).
Recursive Backtracking:     1,540,613ms (~25.68 minutes).
Brute Force:                1,770,851ms (~29.51 minutes).
Short-Circuited (Basic):    49,874ms (~50 seconds).
Parallel:                   10,589ms (~10.5 seconds).

On all 100 equations in Flat200-479:

Clause VSIDS:               19,181ms (~19.18 seconds).
M-size (10\%):              17,716ms (~17.72 seconds).
AMA:                        24,238ms (~24.24 seconds).

On all 1000 equations in UUF50.218.1000:

Clause VSIDS:               1,804ms (~1.80 seconds).
M-size (10\%)               1,768ms (~1.77 seconds).

AMA:                        1,629ms (~1.63 seconds).
EMA (X = 1, Y = 0.95):      1,504ms (~1.50 seconds).

Improved CDCL:              2,385ms (~2.39 seconds).
First UIPs:                 14,012ms (~14.01 seconds).
Backjumping:                11,291ms (~11.29 seconds).
Clause Learning:            9,652ms (~9.65 seconds).
Iterative DPLL (AbP):       6,369ms (~6.37 seconds).
Iterative DPLL (PbA):       6,365ms (~6.37 seconds).
Watched Literals:           7,130ms (~7.13 seconds).
Pure Literal Elimination:   231,107ms (~3.85 minutes).
Unit Propagation:           140,041ms (~2.34 minutes).
\end{FVerbatim}
\noindent
Sorry that this section is really bloated.

\section{Restarts}

This is the last thing we will implement in our SAT solver (for now, maybe). 
Restarts are exactly what it sounds like. We completely delete the trail and 
all notion of assignment. Why would we do that? A valid question. The idea is that 
perhaps while we've been solving, we got stuck in a bad part of the search space
and are just getting contradiction after contradiction. Rather than staying stuck there,
perhaps we could restart our solver so that it will end up in a different part of the 
search space that might be better. 

Makes enough sense, but how do we ensure that we don't just go back to the same 
part of the search space again? Well luckily because we've been learning, our solver 
isn't exactly the same as it initially was. It ``knows'' more, and since knowledge is 
power, the new learned clauses can force certain assignments and propagations earlier 
and that will nudge or push us into the right direction. 

So implement this, it's really really simple. All we need to do is after a certain number 
of conflicts, clear ONLY the trail. We don't clear out our VSIDS, clause measurements, 
reset the watched literals, nothing, nada, nahin. Why? Because the idea of restarts is 
that when the going gets tough, we get going. Notice how I said that we get going, not 
that we get forgetful? We go back to the beginning, remembering everything we've learnt, 
and use that to (hopefully) not repeat mistakes. If we forgot things like our VSIDS, we 
would make the same assignment with the exact same decision, and we'd have to rely only 
on our learned clauses to sway our decision making. But since VSIDS tells us which 
variables have been important in conflicts, it tells us which variables are best to start 
out with on our new try.

For our solver, I chose to restart every time we clear out the clause database. If you've 
been loosy-goosy about following my implementation or have been doing your own experiments
(nice), you can define your own condition. Just remember these things when you're doing a 
restart:
\begin{itemize}
    \item Clear the entire trail.
    \item The trail\_starts array must always have a decision level of 0.
    \item Clearing out var\_to\_trail is not strictly necessary as long 
            as your implementation is correct.
    \item If you do clear out var\_to\_trail, make sure it's resized properly.
    \item Reset qHead.
    \item Reset the decision level to 0.
    \item Unassign every variable.
    \item You might want to do a root level propagation.
\end{itemize}

Here are the results:
\begin{FVerbatim}
On all 1000 equations in uf20-91:

Restarts:                   385ms (~0.39 seconds).
Clause VSIDS:               347ms (~0.35 seconds).
M-size (10\%):              320ms (~0.32 seconds).

AMA (Y = 0.75 or 0.99):     274ms (~0.27 seconds).
EMA (X = 1, Y = 0.95):      266ms (~0.27 seconds).

Improved CDCL:              284ms (~0.28 seconds).
First UIPs:                 406ms (~0.41 seconds).
Backjumping:                380ms (~0.38 seconds).
Clause Learning:            425ms (~0.43 seconds).
Iterative DPLL (AbP):       328ms (~0.33 seconds).
Iterative DPLL (PbA):       324ms (~0.32 seconds).
Watched Literals:           349ms (~0.35 seconds).
Pure Literal Elimination:   1548ms (~1.55 seconds).
Unit Propagation:           868ms (~0.87 seconds).
Recursive Backtracking:     1,540,613ms (~25.68 minutes).
Brute Force:                1,770,851ms (~29.51 minutes).
Short-Circuited (Basic):    49,874ms (~50 seconds).
Parallel:                   10,589ms (~10.5 seconds).

On all 100 equations in Flat200-479:

Restarts:                   9,428ms  (~9.43 seconds).
Clause VSIDS:               19,181ms (~19.18 seconds).
M-size (10\%):              17,716ms (~17.72 seconds).
AMA:                        24,238ms (~24.24 seconds).

On all 1000 equations in UUF50.218.1000:

Restarts:                   1,824ms (~1.82 seconds).
Clause VSIDS:               1,804ms (~1.80 seconds).
M-size (10\%)               1,768ms (~1.77 seconds).

AMA:                        1,629ms (~1.63 seconds).
EMA (X = 1, Y = 0.95):      1,504ms (~1.50 seconds).

Improved CDCL:              2,385ms (~2.39 seconds).
First UIPs:                 14,012ms (~14.01 seconds).
Backjumping:                11,291ms (~11.29 seconds).
Clause Learning:            9,652ms (~9.65 seconds).
Iterative DPLL (AbP):       6,369ms (~6.37 seconds).
Iterative DPLL (PbA):       6,365ms (~6.37 seconds).
Watched Literals:           7,130ms (~7.13 seconds).
Pure Literal Elimination:   231,107ms (~3.85 minutes).
Unit Propagation:           140,041ms (~2.34 minutes).
\end{FVerbatim}
\noindent
Even though our solver is getting slower and slower on uf20-91, I think it's
actually really good that for the larger problems, it's getting substantially 
faster. Even though we want our algorithms and programs to work really quickly 
on small inputs, I think it's important for us to at some point prioritize the 
larger inputs. And given how most SAT problems are even larger than these, 
I think that's a good thing. 

\subsection{Assignment Polarity}

Back in the day, we had this thing called ``Pure Literal Elimination,'' but when 
we transitioned to CDCL, it went away \textit{forever}. Well dear reader, I'm here 
to tell you that it's still going to be gone forever, but its spirit lives on with us.
And by that, I mean, we're going to take the idea behind pure literal elimination and 
apply it to our restarts. 

You see, as we continue learning new clauses, we also encounter new literals. Earlier, 
back when we switched from brute force to recursive backtracking, we noticed that solving 
times decreased. Other than the main structure of the solvers, the only thing that changed
was in brute force, we started out with most variable assignments being False, while in 
recursive backtracking and beyond, they start out as True. I made a claim that this speedup 
was because there were less negative literals than there were positive literals. I won't 
say that I knew for sure, but if we stop for a second to think about it, it makes sense. 

When there are more positive literals, the likelihood of a variable's final assignment 
being True is higher than when there are more negative literals. This was the idea behind 
pure literal elimination, to an extent. Pure literal elimination worked because when a
literal only shows up with one polarity, if there is a solution, there is a 100\% chance 
that there is at least one solution where that pure literal's assignment matches its
polarity (T for positive, F for negative). We want to bring this idea into our assignment 
polarity. While we can't guarantee that literals will remain pure because of clause learning, 
we can at least make better decisions based on the current number of positive/negative occurrences 
of a variable.

To do this, we'll add in a \textit{polarity} array to our solver:
\begin{lstlisting}
    public interface:

    private interface:
        ...
        int[] polarity
        ...
\end{lstlisting}
\noindent
Instead of using a count for both positive and negative literals, instead, we'll 
use one counter for each variable. When the count is negative, we know that there 
are more negative occurrences of the variable, and when the count is nonnegative, we 
will pretend that are more positive occurrences of the variable (even though when the 
count is 0, that means it's balanced).

Because we're using one counter per variable, it might be tempting to use a literal's 
value directly in the count. After all, if we encounter one $+4$ and one $-4$, they'd 
cancel out to 0. Although this is logically correct, I'm worried about integer overflow 
for equations with millions of variables
\footnote{$2^32 \approx 4.3$ billion, so if we had one million variables, we could only 
handle ~43 occurrences of one polarity, which is pretty low, all things considered.}

Instead we'll settle for an ``if negative => count - 1, else count + 1'' kind of deal.
We will gather the counts in the setup phase of our solver, and we will update the counts 
only when we learn a new clause. However, it would be interesting to update the counts for 
all variables involved in a conflict, and see how that affects things. Because of how 
easy this code is to add as well, I'm gonna let you figure it out. Lastly, when we 
pick a variable via its VSIDS, we assign it based on its polarity. Also very easy to 
figure out, you got this, I believe in you.

Here are the results:
\begin{FVerbatim}
On all 1000 equations in uf20-91:

Polarity:                   369ms (~0.37 seconds).
Restarts:                   385ms (~0.39 seconds).
Clause VSIDS:               347ms (~0.35 seconds).
M-size (10\%):              320ms (~0.32 seconds).

AMA (Y = 0.75 or 0.99):     274ms (~0.27 seconds).
EMA (X = 1, Y = 0.95):      266ms (~0.27 seconds).

Improved CDCL:              284ms (~0.28 seconds).
First UIPs:                 406ms (~0.41 seconds).
Backjumping:                380ms (~0.38 seconds).
Clause Learning:            425ms (~0.43 seconds).
Iterative DPLL (AbP):       328ms (~0.33 seconds).
Iterative DPLL (PbA):       324ms (~0.32 seconds).
Watched Literals:           349ms (~0.35 seconds).
Pure Literal Elimination:   1548ms (~1.55 seconds).
Unit Propagation:           868ms (~0.87 seconds).
Recursive Backtracking:     1,540,613ms (~25.68 minutes).
Brute Force:                1,770,851ms (~29.51 minutes).
Short-Circuited (Basic):    49,874ms (~50 seconds).
Parallel:                   10,589ms (~10.5 seconds).

On all 100 equations in Flat200-479:

Polarity:                   8,580ms  (~8.58 seconds).
Restarts:                   9,428ms  (~9.43 seconds).
Clause VSIDS:               19,181ms (~19.18 seconds).
M-size (10\%):              17,716ms (~17.72 seconds).
AMA:                        24,238ms (~24.24 seconds).

On all 1000 equations in UUF50.218.1000:

Polarity:                   1,897ms (~1.90 seconds).
Restarts:                   1,824ms (~1.82 seconds).
Clause VSIDS:               1,804ms (~1.80 seconds).
M-size (10\%)               1,768ms (~1.77 seconds).

AMA:                        1,629ms (~1.63 seconds).
EMA (X = 1, Y = 0.95):      1,504ms (~1.50 seconds).

Improved CDCL:              2,385ms (~2.39 seconds).
First UIPs:                 14,012ms (~14.01 seconds).
Backjumping:                11,291ms (~11.29 seconds).
Clause Learning:            9,652ms (~9.65 seconds).
Iterative DPLL (AbP):       6,369ms (~6.37 seconds).
Iterative DPLL (PbA):       6,365ms (~6.37 seconds).
Watched Literals:           7,130ms (~7.13 seconds).
Pure Literal Elimination:   231,107ms (~3.85 minutes).
Unit Propagation:           140,041ms (~2.34 minutes).
\end{FVerbatim}
\noindent
That's a pretty decent speedup on the larger equations. Not sure how I feel 
about the speedup on UUF50.218.1000, but since it's not even 100ms across all 
1000 cases, I think I'll take it.

\end{document}